% concatenated from main.tex, preambule.tex, math_command.tex

\documentclass{article}

\usepackage{microtype}
\usepackage{graphicx}
\usepackage{subcaption}
\usepackage{booktabs} %

\usepackage{hyperref}


\newcommand{\theHalgorithm}{\arabic{algorithm}}


\usepackage[preprint]{icml2026}


\usepackage{amsmath}
\usepackage{amssymb}
\usepackage{mathtools}
\usepackage{amsthm}
\usepackage{thmtools} %
\usepackage{nicefrac}

\usepackage[capitalize,noabbrev]{cleveref}


\usepackage[disable,textsize=tiny,backgroundcolor=white]{todonotes}




\newcommand{\aymeric}[1]{\todo[bordercolor=Violet, textcolor=Violet, linecolor=Violet, size=\tiny]{\textbf{AD:} #1}}
\newcommand{\hadrien}[1]{\todo[bordercolor=OliveGreen, textcolor=OliveGreen, linecolor=OliveGreen, size=\tiny]{\textbf{HH:} #1}}
\newcommand{\renaud}[1]{\todo[bordercolor=BrickRed, textcolor=BrickRed, linecolor=BrickRed]{\textbf{RG:} #1}}


\usepackage{amsmath,amsfonts,bm}

\usepackage{bbm}

\newcommand{\figleft}{{\em (Left)}}
\newcommand{\figcenter}{{\em (Center)}}
\newcommand{\figright}{{\em (Right)}}
\newcommand{\figtop}{{\em (Top)}}
\newcommand{\figbottom}{{\em (Bottom)}}
\newcommand{\captiona}{{\em (a)}}
\newcommand{\captionb}{{\em (b)}}
\newcommand{\captionc}{{\em (c)}}
\newcommand{\captiond}{{\em (d)}}


\newcommand{\newterm}[1]{{\bf #1}}


\def\figref#1{figure~\ref{#1}}
\def\Figref#1{Figure~\ref{#1}}
\def\twofigref#1#2{figures \ref{#1} and \ref{#2}}
\def\quadfigref#1#2#3#4{figures \ref{#1}, \ref{#2}, \ref{#3} and \ref{#4}}
\def\secref#1{section~\ref{#1}}
\def\Secref#1{Section~\ref{#1}}
\def\twosecrefs#1#2{sections \ref{#1} and \ref{#2}}
\def\secrefs#1#2#3{sections \ref{#1}, \ref{#2} and \ref{#3}}
\def\eqref#1{equation~\ref{#1}}
\def\Eqref#1{Equation~\ref{#1}}
\def\plaineqref#1{\ref{#1}}
\def\chapref#1{chapter~\ref{#1}}
\def\Chapref#1{Chapter~\ref{#1}}
\def\rangechapref#1#2{chapters\ref{#1}--\ref{#2}}
\def\algref#1{algorithm~\ref{#1}}
\def\Algref#1{Algorithm~\ref{#1}}
\def\twoalgref#1#2{algorithms \ref{#1} and \ref{#2}}
\def\Twoalgref#1#2{Algorithms \ref{#1} and \ref{#2}}
\def\partref#1{part~\ref{#1}}
\def\Partref#1{Part~\ref{#1}}
\def\twopartref#1#2{parts \ref{#1} and \ref{#2}}

\def\ceil#1{\lceil #1 \rceil}
\def\floor#1{\lfloor #1 \rfloor}
\def\1{\bm{1}}
\newcommand{\train}{\mathcal{D}}
\newcommand{\valid}{\mathcal{D_{\mathrm{valid}}}}
\newcommand{\test}{\mathcal{D_{\mathrm{test}}}}

\def\eps{{\varepsilon}}


\def\reta{{\textnormal{$\eta$}}}
\def\ra{{\textnormal{a}}}
\def\rb{{\textnormal{b}}}
\def\rc{{\textnormal{c}}}
\def\rd{{\textnormal{d}}}
\def\re{{\textnormal{e}}}
\def\rf{{\textnormal{f}}}
\def\rg{{\textnormal{g}}}
\def\rh{{\textnormal{h}}}
\def\ri{{\textnormal{i}}}
\def\rj{{\textnormal{j}}}
\def\rk{{\textnormal{k}}}
\def\rl{{\textnormal{l}}}
\def\rn{{\textnormal{n}}}
\def\ro{{\textnormal{o}}}
\def\rp{{\textnormal{p}}}
\def\rq{{\textnormal{q}}}
\def\rr{{\textnormal{r}}}
\def\rs{{\textnormal{s}}}
\def\rt{{\textnormal{t}}}
\def\ru{{\textnormal{u}}}
\def\rv{{\textnormal{v}}}
\def\rw{{\textnormal{w}}}
\def\rx{{\textnormal{x}}}
\def\ry{{\textnormal{y}}}
\def\rz{{\textnormal{z}}}

\def\rvepsilon{{\mathbf{\epsilon}}}
\def\rvtheta{{\mathbf{\theta}}}
\def\rva{{\mathbf{a}}}
\def\rvb{{\mathbf{b}}}
\def\rvc{{\mathbf{c}}}
\def\rvd{{\mathbf{d}}}
\def\rve{{\mathbf{e}}}
\def\rvf{{\mathbf{f}}}
\def\rvg{{\mathbf{g}}}
\def\rvh{{\mathbf{h}}}
\def\rvu{{\mathbf{i}}}
\def\rvj{{\mathbf{j}}}
\def\rvk{{\mathbf{k}}}
\def\rvl{{\mathbf{l}}}
\def\rvm{{\mathbf{m}}}
\def\rvn{{\mathbf{n}}}
\def\rvo{{\mathbf{o}}}
\def\rvp{{\mathbf{p}}}
\def\rvq{{\mathbf{q}}}
\def\rvr{{\mathbf{r}}}
\def\rvs{{\mathbf{s}}}
\def\rvt{{\mathbf{t}}}
\def\rvu{{\mathbf{u}}}
\def\rvv{{\mathbf{v}}}
\def\rvw{{\mathbf{w}}}
\def\rvx{{\mathbf{x}}}
\def\rvy{{\mathbf{y}}}
\def\rvz{{\mathbf{z}}}

\def\erva{{\textnormal{a}}}
\def\ervb{{\textnormal{b}}}
\def\ervc{{\textnormal{c}}}
\def\ervd{{\textnormal{d}}}
\def\erve{{\textnormal{e}}}
\def\ervf{{\textnormal{f}}}
\def\ervg{{\textnormal{g}}}
\def\ervh{{\textnormal{h}}}
\def\ervi{{\textnormal{i}}}
\def\ervj{{\textnormal{j}}}
\def\ervk{{\textnormal{k}}}
\def\ervl{{\textnormal{l}}}
\def\ervm{{\textnormal{m}}}
\def\ervn{{\textnormal{n}}}
\def\ervo{{\textnormal{o}}}
\def\ervp{{\textnormal{p}}}
\def\ervq{{\textnormal{q}}}
\def\ervr{{\textnormal{r}}}
\def\ervs{{\textnormal{s}}}
\def\ervt{{\textnormal{t}}}
\def\ervu{{\textnormal{u}}}
\def\ervv{{\textnormal{v}}}
\def\ervw{{\textnormal{w}}}
\def\ervx{{\textnormal{x}}}
\def\ervy{{\textnormal{y}}}
\def\ervz{{\textnormal{z}}}

\def\rmA{{\mathbf{A}}}
\def\rmB{{\mathbf{B}}}
\def\rmC{{\mathbf{C}}}
\def\rmD{{\mathbf{D}}}
\def\rmE{{\mathbf{E}}}
\def\rmF{{\mathbf{F}}}
\def\rmG{{\mathbf{G}}}
\def\rmH{{\mathbf{H}}}
\def\rmI{{\mathbf{I}}}
\def\rmJ{{\mathbf{J}}}
\def\rmK{{\mathbf{K}}}
\def\rmL{{\mathbf{L}}}
\def\rmM{{\mathbf{M}}}
\def\rmN{{\mathbf{N}}}
\def\rmO{{\mathbf{O}}}
\def\rmP{{\mathbf{P}}}
\def\rmQ{{\mathbf{Q}}}
\def\rmR{{\mathbf{R}}}
\def\rmS{{\mathbf{S}}}
\def\rmT{{\mathbf{T}}}
\def\rmU{{\mathbf{U}}}
\def\rmV{{\mathbf{V}}}
\def\rmW{{\mathbf{W}}}
\def\rmX{{\mathbf{X}}}
\def\rmY{{\mathbf{Y}}}
\def\rmZ{{\mathbf{Z}}}

\def\ermA{{\textnormal{A}}}
\def\ermB{{\textnormal{B}}}
\def\ermC{{\textnormal{C}}}
\def\ermD{{\textnormal{D}}}
\def\ermE{{\textnormal{E}}}
\def\ermF{{\textnormal{F}}}
\def\ermG{{\textnormal{G}}}
\def\ermH{{\textnormal{H}}}
\def\ermI{{\textnormal{I}}}
\def\ermJ{{\textnormal{J}}}
\def\ermK{{\textnormal{K}}}
\def\ermL{{\textnormal{L}}}
\def\ermM{{\textnormal{M}}}
\def\ermN{{\textnormal{N}}}
\def\ermO{{\textnormal{O}}}
\def\ermP{{\textnormal{P}}}
\def\ermQ{{\textnormal{Q}}}
\def\ermR{{\textnormal{R}}}
\def\ermS{{\textnormal{S}}}
\def\ermT{{\textnormal{T}}}
\def\ermU{{\textnormal{U}}}
\def\ermV{{\textnormal{V}}}
\def\ermW{{\textnormal{W}}}
\def\ermX{{\textnormal{X}}}
\def\ermY{{\textnormal{Y}}}
\def\ermZ{{\textnormal{Z}}}

\def\vzero{{\bm{0}}}
\def\vone{{\bm{1}}}
\def\vmu{{\bm{\mu}}}
\def\vtheta{{\bm{\theta}}}
\def\vxi{{\bm{\xi}}}
\def\va{{\bm{a}}}
\def\vb{{\bm{b}}}
\def\vc{{\bm{c}}}
\def\vd{{\bm{d}}}
\def\ve{{\bm{e}}}
\def\vf{{\bm{f}}}
\def\vg{{\bm{g}}}
\def\vh{{\bm{h}}}
\def\vi{{\bm{i}}}
\def\vj{{\bm{j}}}
\def\vk{{\bm{k}}}
\def\vl{{\bm{l}}}
\def\vm{{\bm{m}}}
\def\vn{{\bm{n}}}
\def\vo{{\bm{o}}}
\def\vp{{\bm{p}}}
\def\vq{{\bm{q}}}
\def\vr{{\bm{r}}}
\def\vs{{\bm{s}}}
\def\vt{{\bm{t}}}
\def\vu{{\bm{u}}}
\def\vv{{\bm{v}}}
\def\vw{{\bm{w}}}
\def\vx{{\bm{x}}}
\def\vy{{\bm{y}}}
\def\vz{{\bm{z}}}

\def\evalpha{{\alpha}}
\def\evbeta{{\beta}}
\def\evepsilon{{\epsilon}}
\def\evlambda{{\lambda}}
\def\evomega{{\omega}}
\def\evmu{{\mu}}
\def\evpsi{{\psi}}
\def\evsigma{{\sigma}}
\def\evtheta{{\theta}}
\def\eva{{a}}
\def\evb{{b}}
\def\evc{{c}}
\def\evd{{d}}
\def\eve{{e}}
\def\evf{{f}}
\def\evg{{g}}
\def\evh{{h}}
\def\evi{{i}}
\def\evj{{j}}
\def\evk{{k}}
\def\evl{{l}}
\def\evm{{m}}
\def\evn{{n}}
\def\evo{{o}}
\def\evp{{p}}
\def\evq{{q}}
\def\evr{{r}}
\def\evs{{s}}
\def\evt{{t}}
\def\evu{{u}}
\def\evv{{v}}
\def\evw{{w}}
\def\evx{{x}}
\def\evy{{y}}
\def\evz{{z}}

\def\mA{{\bm{A}}}
\def\mB{{\bm{B}}}
\def\mC{{\bm{C}}}
\def\mD{{\bm{D}}}
\def\mE{{\bm{E}}}
\def\mF{{\bm{F}}}
\def\mG{{\bm{G}}}
\def\mH{{\bm{H}}}
\def\mI{{\bm{I}}}
\def\mJ{{\bm{J}}}
\def\mK{{\bm{K}}}
\def\mL{{\bm{L}}}
\def\mM{{\bm{M}}}
\def\mN{{\bm{N}}}
\def\mO{{\bm{O}}}
\def\mP{{\bm{P}}}
\def\mQ{{\bm{Q}}}
\def\mR{{\bm{R}}}
\def\mS{{\bm{S}}}
\def\mT{{\bm{T}}}
\def\mU{{\bm{U}}}
\def\mV{{\bm{V}}}
\def\mW{{\bm{W}}}
\def\mX{{\bm{X}}}
\def\mY{{\bm{Y}}}
\def\mZ{{\bm{Z}}}
\def\mBeta{{\bm{\beta}}}
\def\mPhi{{\bm{\Phi}}}
\def\mLambda{{\bm{\Lambda}}}
\def\mSigma{{\bm{\Sigma}}}

\DeclareMathAlphabet{\mathsfit}{\encodingdefault}{\sfdefault}{m}{sl}
\SetMathAlphabet{\mathsfit}{bold}{\encodingdefault}{\sfdefault}{bx}{n}
\newcommand{\tens}[1]{\bm{\mathsfit{#1}}}
\def\tA{{\tens{A}}}
\def\tB{{\tens{B}}}
\def\tC{{\tens{C}}}
\def\tD{{\tens{D}}}
\def\tE{{\tens{E}}}
\def\tF{{\tens{F}}}
\def\tG{{\tens{G}}}
\def\tH{{\tens{H}}}
\def\tI{{\tens{I}}}
\def\tJ{{\tens{J}}}
\def\tK{{\tens{K}}}
\def\tL{{\tens{L}}}
\def\tM{{\tens{M}}}
\def\tN{{\tens{N}}}
\def\tO{{\tens{O}}}
\def\tP{{\tens{P}}}
\def\tQ{{\tens{Q}}}
\def\tR{{\tens{R}}}
\def\tS{{\tens{S}}}
\def\tT{{\tens{T}}}
\def\tU{{\tens{U}}}
\def\tV{{\tens{V}}}
\def\tW{{\tens{W}}}
\def\tX{{\tens{X}}}
\def\tY{{\tens{Y}}}
\def\tZ{{\tens{Z}}}


\def\gA{{\mathcal{A}}}
\def\gB{{\mathcal{B}}}
\def\gC{{\mathcal{C}}}
\def\gD{{\mathcal{D}}}
\def\gE{{\mathcal{E}}}
\def\gF{{\mathcal{F}}}
\def\gG{{\mathcal{G}}}
\def\gH{{\mathcal{H}}}
\def\gI{{\mathcal{I}}}
\def\gJ{{\mathcal{J}}}
\def\gK{{\mathcal{K}}}
\def\gL{{\mathcal{L}}}
\def\gM{{\mathcal{M}}}
\def\gN{{\mathcal{N}}}
\def\gO{{\mathcal{O}}}
\def\gP{{\mathcal{P}}}
\def\gQ{{\mathcal{Q}}}
\def\gR{{\mathcal{R}}}
\def\gS{{\mathcal{S}}}
\def\gT{{\mathcal{T}}}
\def\gU{{\mathcal{U}}}
\def\gV{{\mathcal{V}}}
\def\gW{{\mathcal{W}}}
\def\gX{{\mathcal{X}}}
\def\gY{{\mathcal{Y}}}
\def\gZ{{\mathcal{Z}}}

\def\sA{{\mathbb{A}}}
\def\sB{{\mathbb{B}}}
\def\sC{{\mathbb{C}}}
\def\sD{{\mathbb{D}}}
\def\sF{{\mathbb{F}}}
\def\sG{{\mathbb{G}}}
\def\sH{{\mathbb{H}}}
\def\sI{{\mathbb{I}}}
\def\sJ{{\mathbb{J}}}
\def\sK{{\mathbb{K}}}
\def\sL{{\mathbb{L}}}
\def\sM{{\mathbb{M}}}
\def\sN{{\mathbb{N}}}
\def\sO{{\mathbb{O}}}
\def\sP{{\mathbb{P}}}
\def\sQ{{\mathbb{Q}}}
\def\sR{{\mathbb{R}}}
\def\sS{{\mathbb{S}}}
\def\sT{{\mathbb{T}}}
\def\sU{{\mathbb{U}}}
\def\sV{{\mathbb{V}}}
\def\sW{{\mathbb{W}}}
\def\sX{{\mathbb{X}}}
\def\sY{{\mathbb{Y}}}
\def\sZ{{\mathbb{Z}}}

\def\emLambda{{\Lambda}}
\def\emA{{A}}
\def\emB{{B}}
\def\emC{{C}}
\def\emD{{D}}
\def\emE{{E}}
\def\emF{{F}}
\def\emG{{G}}
\def\emH{{H}}
\def\emI{{I}}
\def\emJ{{J}}
\def\emK{{K}}
\def\emL{{L}}
\def\emM{{M}}
\def\emN{{N}}
\def\emO{{O}}
\def\emP{{P}}
\def\emQ{{Q}}
\def\emR{{R}}
\def\emS{{S}}
\def\emT{{T}}
\def\emU{{U}}
\def\emV{{V}}
\def\emW{{W}}
\def\emX{{X}}
\def\emY{{Y}}
\def\emZ{{Z}}
\def\emSigma{{\Sigma}}

\newcommand{\etens}[1]{\mathsfit{#1}}
\def\etLambda{{\etens{\Lambda}}}
\def\etA{{\etens{A}}}
\def\etB{{\etens{B}}}
\def\etC{{\etens{C}}}
\def\etD{{\etens{D}}}
\def\etE{{\etens{E}}}
\def\etF{{\etens{F}}}
\def\etG{{\etens{G}}}
\def\etH{{\etens{H}}}
\def\etI{{\etens{I}}}
\def\etJ{{\etens{J}}}
\def\etK{{\etens{K}}}
\def\etL{{\etens{L}}}
\def\etM{{\etens{M}}}
\def\etN{{\etens{N}}}
\def\etO{{\etens{O}}}
\def\etP{{\etens{P}}}
\def\etQ{{\etens{Q}}}
\def\etR{{\etens{R}}}
\def\etS{{\etens{S}}}
\def\etT{{\etens{T}}}
\def\etU{{\etens{U}}}
\def\etV{{\etens{V}}}
\def\etW{{\etens{W}}}
\def\etX{{\etens{X}}}
\def\etY{{\etens{Y}}}
\def\etZ{{\etens{Z}}}

\newcommand{\pdata}{p_{\rm{data}}}
\newcommand{\ptrain}{\hat{p}_{\rm{data}}}
\newcommand{\Ptrain}{\hat{P}_{\rm{data}}}
\newcommand{\pmodel}{p_{\rm{model}}}
\newcommand{\Pmodel}{P_{\rm{model}}}
\newcommand{\ptildemodel}{\tilde{p}_{\rm{model}}}
\newcommand{\pencode}{p_{\rm{encoder}}}
\newcommand{\pdecode}{p_{\rm{decoder}}}
\newcommand{\precons}{p_{\rm{reconstruct}}}

\newcommand{\laplace}{\mathrm{Laplace}} %

\newcommand{\E}{\mathbb{E}}
\newcommand{\Ls}{\mathcal{L}}
\newcommand{\R}{\mathbb{R}}
\newcommand{\emp}{\tilde{p}}
\newcommand{\lr}{\alpha}
\newcommand{\reg}{\lambda}
\newcommand{\rect}{\mathrm{rectifier}}
\newcommand{\softmax}{\mathrm{softmax}}
\newcommand{\sigmoid}{\sigma}
\newcommand{\softplus}{\zeta}
\newcommand{\KL}{D_{\mathrm{KL}}}
\newcommand{\Var}{\mathrm{Var}}
\newcommand{\standarderror}{\mathrm{SE}}
\newcommand{\Cov}{\mathrm{Cov}}
\newcommand{\normlzero}{L^0}
\newcommand{\normlone}{L^1}
\newcommand{\normltwo}{L^2}
\newcommand{\normlp}{L^p}
\newcommand{\normmax}{L^\infty}

\newcommand{\parents}{Pa} %

\DeclareMathOperator*{\argmax}{arg\,max}
\DeclareMathOperator*{\argmin}{arg\,min}

\DeclareMathOperator{\sign}{sign}
\DeclareMathOperator{\Tr}{Tr}
\let\ab\allowbreak


\newcommand{\clip}{\mathrm{clip}}
\renewcommand{\tL}{\tilde{L}}
\newcommand{\tmu}{\tilde{\mu}}


\newcommand{\cA}{\mathcal{A}}
\newcommand{\cB}{\mathcal{B}}
\newcommand{\cC}{\mathcal{C}}
\newcommand{\cD}{\mathcal{D}}
\newcommand{\cE}{\mathcal{E}}
\newcommand{\cF}{\mathcal{F}}
\newcommand{\cG}{\mathcal{G}}
\newcommand{\cH}{\mathcal{H}}
\newcommand{\cI}{\mathcal{I}}
\newcommand{\cJ}{\mathcal{J}}
\newcommand{\cK}{\mathcal{K}}
\newcommand{\cL}{\mathcal{L}}
\newcommand{\cM}{\mathcal{M}}
\newcommand{\cN}{\mathcal{N}}
\newcommand{\cO}{\mathcal{O}}
\newcommand{\cP}{\mathcal{P}}
\newcommand{\cQ}{\mathcal{Q}}
\newcommand{\cR}{\mathcal{R}}
\newcommand{\cS}{\mathcal{S}}
\newcommand{\cT}{\mathcal{T}}
\newcommand{\cU}{\mathcal{U}}
\newcommand{\cV}{\mathcal{V}}
\newcommand{\cW}{\mathcal{W}}
\newcommand{\cX}{\mathcal{X}}
\newcommand{\cY}{\mathcal{Y}}
\newcommand{\cZ}{\mathcal{Z}}




\theoremstyle{plain}
\newtheorem{theorem}{Theorem}[section]
\newtheorem{proposition}[theorem]{Proposition}
\newtheorem{lemma}[theorem]{Lemma}
\newtheorem{corollary}[theorem]{Corollary}
\theoremstyle{definition}
\newtheorem{definition}[theorem]{Definition}
\newtheorem{assumption}[theorem]{Assumption}
\theoremstyle{remark}
\newtheorem{remark}[theorem]{Remark}



\icmltitlerunning{Byzantine Distributed Learning under similarity}

\begin{document}

\twocolumn[
  \icmltitle{From Inexact Gradients to Byzantine Robustness: \\ Acceleration and Optimization under Similarity}
  






  \begin{icmlauthorlist}
    \icmlauthor{Renaud Gaucher}{xxx,inria}
    \icmlauthor{Aymeric Dieuleveut}{xxx}
    \icmlauthor{Hadrien Hendrikx}{inria}
  \end{icmlauthorlist}

  \icmlaffiliation{xxx}{Centre de Mathématiques Appliquées,
    École polytechnique, Institut Polytechnique de Paris, Palaiseau France}
  \icmlaffiliation{inria}{Centre Inria de l'Univ. Grenoble Alpes, CNRS, LJK, Grenoble, France}

  \icmlcorrespondingauthor{Renaud Gaucher}{renaud.gaucher@polytechnique.edu}

  \icmlkeywords{Machine Learning, ICML, Distributed Optimization, Byzantine Fault}

  \vskip 0.3in
]



\printAffiliationsAndNotice{}  %

\begin{abstract}
Standard federated learning algorithms are vulnerable to adversarial nodes, a.k.a. Byzantine failures. To solve this issue, robust distributed learning algorithms have been developed, which typically replace parameter averaging by robust aggregations. While generic conditions on these aggregations exist to guarantee the convergence of (Stochastic) Gradient Descent (SGD), the analyses remain rather ad-hoc. This hinders the development of more complex robust algorithms, such as accelerated ones.
In this work, we show that Byzantine-robust distributed optimization can, under standard generic assumptions, be cast as a general optimization with inexact gradient oracles (with both additive and multiplicative error terms), an active field of research. 

This allows for instance to directly show that GD on top of standard robust aggregation procedures obtains optimal asymptotic error in the Byzantine setting.
Going further, we propose two optimization schemes to speed up the convergence. The first one is a Nesterov-type accelerated scheme whose proof directly derives from accelerated inexact gradient results applied to our formulation. 
The second one hinges on \emph{Optimization under Similarity}, in which the server leverages an auxiliary loss function that  approximates the global loss.
Both approaches allow to drastically reduce the communication complexity compared to previous methods, as we show theoretically and empirically.
\end{abstract}


\section{Introduction}


Distributed Learning is a paradigm in which multiple computers collaborate to train a machine learning model. This allows to simultaneously leverage the computational power of multiple devices and to adapt to the distributed nature of data, which is often stored on smartphones, IOT devices, or by larger stakeholders such as companies or hospitals. However, involving multiple parties in an algorithmic process opens up to new threats: machines can crash, software can be buggy or even hacked, data can be corrupted or even adversarially crafted. Byzantine failure \citep{lamport1982byzantine} is a worst-case model of such behaviors, corresponding to omniscient adversarial parties. 

Distributed Gradient Descent (DGD), the workhorse of modern distributed learning, is vulnerable to such failures \citep{blanchard2017machine}: when the clients' gradients are averaged, one Byzantine node returning adversarially crafted gradients can corrupt the result arbitrarily, thus ruining the optimization. To address this critical issue, a large body of research has been devoted to developing and studying distributed learning methods that converge even with adversaries \citep{blanchard2017machine, yin2018byzantine, chen2017distributed}. Those methods crucially replace the gradients (or local parameters) averaging step by a robust aggregation scheme, such as the geometric median or coordinate-wise trimming. Over the years, analyses specific to each robust aggregation method have given way to generic characterizations and standard study assumptions \citep{allouah2023fixing, karimireddy2020byzantine}. Typically, one simply needs to guarantee that the error (\emph{i.e.}, distance from the robust aggregation result to the average of the honest nodes) remains bounded. \renaud{Je trouve ces deux phrases un peu obscures}This bound depends on both the heterogeneity between honest nodes, and the proportion of Byzantine nodes. Controlling these two sources of error, one can determine the optimal asymptotic error, and show that it can be achieved by several robust aggregation schemes \citep{farhadkhani2022byzantine, karimireddy2021learning}. %
Momentum-based methods were then proposed to alleviate the stochastic error in the Byzantine context \citep{karimireddy2021learning, farhadkhani2022byzantine}. Similarly, refined  results taking precisely the dimension into account have been reached by leveraging high-dimensional averaging tools \citep{diakonikolas2019robust, allouah2025towards}. Robust distributed learning algorithms were also coupled with gradient compression \citep{rammal2024communication, malinovsky2024byzantine, liu2024byzantine} to reduce their communication cost.

As mentioned above, Byzantine-resilient algorithms work for a wide variety of robust aggregation methods, as long as they satisfy relevant error conditions. Yet, analyses of robust algorithms are still tailored to the Byzantine setting: proofs closely follow existing schemes, but are adapted to match the specificity of the distributed adversarial setting. Although this adaptation is rather transparent in simple settings, it becomes much more cumbersome for more involved proofs, with less leeway to accumulate errors terms. This hinders the development of efficient algorithms, leaving critical questions open, and preventing the adaptation of efficient federated algorithms. For instance, \citet{ghosh2020distributed} proposed second order methods, but require a very low heterogeneity level to actually converge. Similarly, \citet{xu2025nesterov} studied Nesterov's acceleration in a relatively homogeneous setting, and do not actually obtain expected accelerated convergence rates. %





 





\subsection*{Contributions and outline}

Our work alleviates the aforementioned shortcomings of existing theory. We cast Byzantine-robust optimization as a special case of optimization under inexact (gradient) oracles, and build on this to derive efficient robust algorithms with provably low iteration (and so communication) complexity.  

\textbf{Generic Formalism.} After reviewing the setting in \Cref{subsec:setup},  we show in \Cref{subsec:reduction} that Byzantine-robust distributed optimization exactly fits into the formalism of first-order methods with inexact gradient oracles. To do so,
we show that combining standard generic assumptions on the heterogeneity of loss functions with a standard characterization of the employed robust aggregation method leads to simple additive and multiplicative error term on the mean of the honest gradient.
Optimization under the provided inexact gradient formulation is an active field of research, from which we can thus directly reuse most results. 

We then show in \Cref{subsec:tightness} that this abstraction is tight in the sense that using existing inexact GD results allows to match the Byzantine-robust optimization lower bound on the asymptotic error under heterogeneity. %
This allows for a more systematic search of efficient algorithms, abstracting away the complexity of the Byzantine setting, and 
thus building on the inexact gradient formulation to propose a robust version of two standard fast algorithms.

\textbf{Acceleration.}  In \Cref{sec:accelerated}, we exhibit a Nesterov-type accelerated robust distributed algorithm which, under medium heterogeneity assumptions, converges with an accelerated linear rate of convergence $\mathcal{O}(\sqrt{\nicefrac{\mu}{L}})$, where $\mu$ is the strong convexity of the considered loss function and $L$ the smoothness. This accelerated rate is obtained as a direct consequence of a result on acceleration under inexact gradients~\citep{devolder2014first}, which illustrates the effectiveness and modularity of our approach.

\textbf{Optimization under similarity.} Next,  if an approximation (in the second-order sense) of the global honest loss is available, we give in \Cref{sec:similarity} a Byzantine-robust algorithm (with computationally more expensive updates than GD) called Prox Inexact Gradient method under Similarity (PIGS). PIGS enjoys a linear convergence rate of $\mathcal{O}(\nicefrac{\Delta}{\mu}\log(1/\eps))$ to reach an $\eps$ neighborhood of the optimal error, where $\Delta$ is a measure of the distance between the Hessian of approximate loss and the true honest one. Since $\Delta$ can be drastically smaller than the smoothness of the loss function, this results in an important gain in iteration complexity, which consequently reduces the communication time. The convergence theory of PIGS requires an extension of~\citet{woodworth2023two} to handle the inexactness of the error term.

\textbf{Experiments.} Last, we empirically show the robustness and performance of the proposed methods by comparing them to existing methods in \Cref{sec:experiments}.


\paragraph{Notations.} We denote $\langle \cdot, \cdot \rangle$ the standard dot product in $\R^d$,  $\|\cdot\|$ the euclidean norm,  and $\|\cdot\|_{op}$ the associated operator norm in $\R^{d\times d}$. For a real valued function $f:\R^d \rightarrow \R$, $\nabla f$ denotes its gradient, and $\nabla^2 f$ its hessian. Let $f(x) \in \Omega(g(x))$ and $g(x) \in \cO(f(x))$ denote that $\exists M>0$ such that $f(x) \ge M g(x)$. We denote $f(x) \in \Theta(g(x))$ when both $f(x)\in \cO(g(x))$ and $f(x) \in \Theta(g(x))$ hold.

\section{Setting and Reduction}

\subsection{Byzantine Robust Distributed Optimization}
\label{subsec:setup}
We consider a server-client distributed system, with $n$ clients. Among them, an unknown subset $\cB \subset [n]$ is composed of $|\cB| = f$ Byzantine clients. We denote $\cH = [n]\backslash \cB$ the $n-f$ honest clients. We aim at finding 
\begin{equation}
\label{eq:byzantine_loss}
    \vx^* =\argmin_{\vx\in\R^d}\left \{ \cL_{\cH}(\vx) := \frac{1}{n-f}\sum_{i \in \cH}\cL_i(\vx)\right\},
\end{equation}
where $\cL_i : \R^d \rightarrow \R$ denotes the loss function of the client~$i$. This can for instance be an empirical risk minimization problem, where $\cL_i$ denotes the empirical loss of the machine learning model on the dataset stored by client $i$. We further assume that all local losses $\cL_i$ are differentiable, convex, and $L$-smooth, and that $\cL_{\cH}$ is $\mu$-strongly convex.
\begin{definition}
\label{def:strong_convexity_and_smoothness}
A function $\cL$ is $\mu$\textit{-strongly convex}, if
    \[
    \forall \vx,\vy \in \R^d, \cL(\vx) \ge \cL(\vy) + \langle  \nabla \cL(\vy), \vx-\vy\rangle + \frac{\mu}{2}\|\vx - \vy\|^2,
    \]
and only \textit{convex} if the above inequality holds with $\mu=0$.

A function $\cL$ is $L$\textit{-smooth} if the gradient $L$-Lipschitz, i.e.
$$\forall \vx,\vy\in \R^d, \|\nabla \cL(\vx) - \nabla \cL(\vy)\|^2 \le L\|\vx-\vy\|^2.$$
The ratio $\kappa = L/\mu$ is called the \emph{condition number} of $\cL$.
\end{definition}

The goal of (distributed) optimization is to output a model $\hat{\vx}$ with an error that vanishes with the computing budget. We say that an algorithm is resilient if it is able to reach a given error level $\eps$ despite the presence of adversaries, namely that the output $\hat{\vx}$ satisfies
\[
\cL_{\cH}(\hat{\vx}) - \cL_{\cH}(\vx^*) \le \eps.
\]


\textbf{Hardness of Byzantine Optimization.}
When honest clients have heterogeneous loss functions, it is generally impossible to achieve arbitrarily precise solutions of \cref{eq:byzantine_loss} in the presence of adversaries, as Byzantine clients cannot be distinguished from honest clients with outlying losses. Consequently, the error increases with both the heterogeneity among honest clients and the proportion of Byzantine clients, up to a critical fraction known as the\textit{ breakdown point}. The following assumption captures these parameters.

\begin{assumption}[(G,B)-heterogeneity]
\label{asmpt:GB_heterogeneity}
    The local loss functions of honest clients are said to satisfy {$(G,B)$-heterogeneity} when, for any $\vx \in \R^d$, 
    \[
    \frac{1}{|\cH|}\sum_{i \in \cH} \|\nabla \cL_i(\vx) - \nabla \cL_{\cH}(\vx)\|^2 \le G^2 + B^2\|\nabla \cL_{\cH}(\vx)\|^2.
    \]
\end{assumption}

    


\begin{theorem}[\citet{allouah2024robust}]
\label{thm:allouha_lower_bound}
  For any distributed algorithm, there exist quadratic local loss functions satisfying \cref{asmpt:GB_heterogeneity}, with an $L$-smooth and $\mu$-strongly convex global loss $\cL_{\cH}$, for which  algorithm's  output $\hat{\vx}$ can have a non-vacuous guarantee only if $ \frac{f}{n}\le \frac{1}{B^2 + 2},$ and  such that
    \begin{align}
     \cL_{\cH}(\hat{\vx}) - \cL_{\cH}(\vx^*) &\ge \frac{G^2}{8\mu} \, \frac{f}{n - (2+B^2)f}. \label{eq:lb_functionvalue}\\
 \|\nabla \cL_{\cH}(\hat{\vx})\|^2 &\ge \frac{\:G^2}{4}\,\frac{f}{n - (2+B^2)f}. \label{eq:lb_gradnorm}
\end{align}
\end{theorem}
The  inequality $ \frac{f}{n}\le \frac{1}{B^2 + 2},$ gives an upper bound on the breakdown point. The two subsequent inequalities lower bounds the reachable target precision (provided that the first one holds), respectively in terms of function values and squared gradient norm.


\subsection{Reduction to an Inexact Gradient Problem}
\label{subsec:reduction}
To alleviate the influence of adversaries, distributed optimization methods rely on estimating robustly the average of the honest gradients, by using for instance a geometric median \citep{small1990survey, pillutla2022robust}, a coordinate-wise trimmed mean \citep{yin2018byzantine}, or Krum \citep{blanchard2017machine}. \citet{allouah2023fixing} proposed a criterion to simultaneously analyze those robust averaging tools, namely by showing that they all are $(f,\nu)$-robust aggregation rules.


\begin{definition}[$(f,\nu)$-robustness \citep{allouah2023fixing}]
\label{def:f_kappa_robustness}
    An aggregation rule $F:\R^{d\times n}\rightarrow \R^d$ is said to be $(f,\nu)$-robust, if $\,\forall \vx_1,\ldots,\vx_n\in \R^d$, and any set $S\subset [n]$ of size $n-f$, 
    \[
    \|F(\vx_1,\ldots,\vx_n) - \overline{\vx}_S\|^2 \le \nu \cdot \frac{1}{|S|}\sum_{i\in S}\|\vx_i - \overline{\vx}_S\|^2,
    \]
    where $\overline{\vx}_S = |S|^{-1}\sum_{i\in S}\vx_i$. We call $\nu$\footnote{We denote it $\nu$ instead of $\kappa$ as in \citet{allouah2023fixing} due to notation incompatibility with the standard notation of the condition number.} the \textit{robustness coefficient} of $F$.
\end{definition}

Closed-form robustness coefficients can be found in \citet{allouah2023fixing}. 
We recall them in \cref{app:sec:robust_aggregation} and complete by extending their mixing strategies. 



The strength of \cref{def:f_kappa_robustness} is  its compatibility with \cref{asmpt:GB_heterogeneity}. 
\begin{lemma}
\label{lemma:Byz_to_Inexact_Gradients}
    Assume that all clients declare a gradient $\nabla \cL_i(\vx)$ to the server, where honest nodes declare their true gradient, while Byzantine nodes declare some arbitrary value. Then, computing $\tilde{\nabla}\cL_{\cH}(\vx) := F(\nabla\cL_1(\vx),\ldots, \nabla \cL_n(\vx))$ yields, if $F$ is $(f,\nu)$-robust and local loss are $(G,B)$-heterogeneous,
\begin{equation}
    \|\tilde{\nabla}\cL_{\cH}(\vx) - \nabla \cL_{\cH}(\vx)\|^2 \le \nu G^2 + \nu B^2 \|\nabla \cL_{\cH}(\vx)\|^2.
\end{equation}
\end{lemma}

\begin{algorithm}[t]
  \caption{Inexact Oracles Sampler}
  \label{alg:reduction}
  \begin{algorithmic}
    \STATE {\bfseries Input:} $F$, $\vx$
    \STATE Server sends $\vx$ to clients
    \FOR{client $i\in[n]$ in parallel} 
    \STATE Computes $
        \vg_i = \begin{cases}
            \nabla \cL_i(\vx) \text{ if } i\in \cH\\
            * \text{ if } i \in \cB
        \end{cases}
    $
    \STATE Send $\vg_i$ to Server
    \ENDFOR
    \STATE Server computes $\tilde \nabla \cL_{\cH}(\vx) := F(\vg_1,\ldots, \vg_n)$
  \end{algorithmic}
\end{algorithm}


In other words, robustly averaging the gradients sent by the honest clients produces a $(\zeta^2,\alpha)$-inexact gradient oracle, as defined hereafter, with $\zeta^2 = \nu G^2$ and $\alpha = \nu B^2$. 
\begin{definition}[$(\zeta^2,\alpha)$-inexact oracle]
\label{def:inexact_oracle}
    A $(\zeta^2,\alpha)$-inexact gradient oracle of function $h$, is an oracle $\tilde \nabla h$ satisfying, for any $\vx \in \R^d$,
    \[
    \|\tilde \nabla h(\vx) - \nabla h(\vx)\|^2 \le \zeta^2 + \alpha \|\nabla h(\vx)\|^2.
    \]
\end{definition}

This reduction is of interest since instances of $(\zeta^2,\alpha)$-inexact oracle have been widely investigated recently \citep{ajalloeian2020convergence}, while mostly with either $\alpha=0$ or $\zeta^2 =0$ \citep{de2020worst, cyrus2018robust, vasin2023accelerated}. We can then reuse these algorithms to obtain Byzantine-robust methods. Moreover, developing algorithms based on the $(\zeta^2,\alpha)$ inexact gradient assumption opens up to generic contributions, of interest beyond the Byzantine setting. 

\subsection{Tightness of the reduction}
\label{subsec:tightness}
A natural question is that of the tightness of the reduction. We hereafter establish that, in terms of the achievable error,  we do not  lose anything by analyzing the Byzantine setup under \Cref{asmpt:GB_heterogeneity} as an instance of a black-box model with a $(\zeta^2,\alpha)$-inexact gradient oracle. 


\begin{proposition}
Reducing \cref{asmpt:GB_heterogeneity} and \cref{def:f_kappa_robustness} to \cref{def:inexact_oracle}, with  $\zeta^2 = \nu G^2$ and $\alpha = \nu B^2$ is tight, in the sense that,  for an optimal $\nu$ in \cref{def:f_kappa_robustness},  the upper bound achieved by GD under the reduced assumption matches the lower bound under \cref{asmpt:GB_heterogeneity} given by \cref{thm:allouha_lower_bound}.
\end{proposition} 

Indeed, we first leverage the guarantee by \citet{ajalloeian2020convergence}: for an $L$-smooth and $\mu$-strongly convex function, gradient descent based on $(\zeta^2,\alpha)$-inexact oracles,  with fixed step size $\eta\le L^{-1}$, is guaranteed to converge to a $\nicefrac{\zeta^2}{2\mu(1-\alpha)}$-neighborhood of the minimal function value. Second, as shown in \citet{allouah2023fixing}, a $(f,\nu)$-robust aggregation rules must satisfy $\nu \ge \nicefrac{f}{n-2f}$ and there exists aggregation rules with $\nu = \cO(\nicefrac{f}{n-f})$ (see \cref{app:sec:robust_aggregation}). Putting things together, with our reduction, $\alpha = \nu B^2= O(\frac{f}{n-2f}B^2)$ and $\zeta^2 = \nu G^2= O(\frac{f}{n-2f}G^2)$, then the asymptotic error is $\cO(\frac{G^2}{\mu}\frac{f }{ n-(2+B^2) f})$, which matches (up to constant) the lower bound (\ref{eq:lb_functionvalue}) of \Cref{thm:allouha_lower_bound}.


This result is somehow surprising, as \cref{alg:reduction} never takes into account that the gradients $(g_i)_{i\in[n]}$ are associated with specific losses, as $(f,\nu)$-robust aggregation rules are usually permutation invariant. This shows that we can attain the optimization lower bound  without observing the identities of the $g_i$.







\begin{remark}
    We can reinterpret the lower bound (\ref{eq:lb_gradnorm}) in \cref{thm:allouha_lower_bound} as describing the region of $\vx \in \R^d$ on which no information can be leveraged by the server. Indeed, 
\cref{eq:lb_gradnorm} gives the gradient  norm value below which 
a $(\zeta^2=\nu G^2, \alpha = \nu B^2)$-inexact oracle resulting from \cref{alg:reduction} may not be informative,
as if $\|\nabla\cL_{\cH}(\vx)\|^2 \le \frac{G^2 \nu}{1 - B^2 \nu}$, then the value $\tilde{\nabla}\cL_{\cH}(\vx) =0$  is a valid $(\zeta^2=\nu G^2, \alpha = \nu B^2)$-inexact oracle, while being obviously un-informative.
\end{remark}







\section{Accelerated Inexact Gradient Descent}
\label{sec:accelerated}

In the following section, we build on an existing notion of inexact gradient oracles~\citep{devolder2013first, devolder2014first} to provide the first accelerated convergence rate for a Byzantine-resilient first-order method.

\subsection{Link to First Order $(\delta,L,\mu)$ Oracles}

The study of inexact gradient oracles dates back at least to \citet{d2008smooth}. \citet{devolder2013first}  introduce the notion of $(\delta, L,\mu)$-oracle of a function $f$. 
\begin{definition}
    $(f_{\delta,L,\mu}, g_{\delta,L,\mu})$ are $(\delta, L,\mu)$-oracle of a convex function $f$ if, for any $\vx,\vy \in \R^d$, it satisfies 
    \begin{align*}
        \frac{\mu}{2}\|\vx-\vy\|^2 \le f(\vx) - &f_{\delta, L,\mu}(\vy) - \langle g_{\delta, L,\mu}(\vy), \vx-\vy \rangle \\ &\le \frac{L}{2}\|\vx-\vy\|^2 + \delta.
    \end{align*}
\end{definition}

The class of $(\delta, L,\mu)$ oracles covers a wide range of scenarios: they appear when the gradients are evaluated at incorrect parameters, when gradients are only Hölder-continuous rather than $L$-Lipschitz, when the oracle relies on an optimization sub-routine, or when gradients come from of a smoothed version of the original loss \citep{devolder2013first, devolder2014first}. Closer to our interest, $(\zeta^2,\alpha)$-inexact gradient oracles are also $(\delta,L,\mu)$ oracles when~${\alpha = 0}$.



\begin{proposition}[\citet{devolder2014first} Section 2.3]
\label{prop:link_beween_inexact_oracles}
 \textbf{1.} Assume that $(f_{\delta, L, \mu},g_{\delta,L,\mu} )$ is a $(\mu, L, \delta)$-oracle of $f$, and that $f$ is continuous and $\mu_f$ strongly convex. Then for $\vx\in\R^d$,
\begin{align*}
        & 0 \le f(\vx) - f_{\delta, L, \mu}(\vx) \le \delta,\\
        & \|\nabla f(\vx) - g_{\delta,L,\mu}(\vx)\|^2 \le   2(L-\mu_f)\delta.
\end{align*}
\textbf{2.} Conversely, if $\|\nabla f(\vx) - g(\vx)\|^2 \le \zeta^2$ and $f$ is $L_f$-smooth and $\mu_f$ strongly convex, then $(f - \nicefrac{\zeta^2}{\mu_f}, g)$ is an $\big((\nicefrac{1}{2L_f}+\nicefrac{1}{\mu_f})\zeta^2, 2L_f,\nicefrac{\mu_f}{2}\big)$-oracle of $(f,\nabla f)$.
\end{proposition}
The above proposition shows that for strongly convex and smooth losses, using the $(\delta, L,\mu)$-oracle formulation is equivalent -- up to a factor $\kappa$ in the inexactness -- to the $(\zeta^2,\alpha=0)$ inexact gradient definition. This allows us to ensure the feasibility of first-order acceleration in distributed optimization under Byzantine corruption.


\subsection{Convergence Results}

\begin{algorithm}[h]
  \caption{Byzantine-Resilient Fast Gradient Method}
  \label{alg:fast_gradient_method}
  \begin{algorithmic}
    \STATE {\bfseries Input:} $F$, $(\gamma_k)$, $(\Gamma_k)$, $\vx_0=\vy_0=\vz_0$
     \FOR{$k=1$ {\bfseries to} $K$}
     \STATE Sample $\tilde\nabla \cL_{\cH}(\vx_k)$ using \cref{alg:reduction}
    \STATE Compute $\vy_{k} = \vx_k - \frac{1}{2L}\tilde{\nabla}\cL_{\cH}(\vx_k)$ 
    \STATE Compute $\vx_{k+1}=\vy_k + \frac{\Gamma_{k-1}}{\Gamma_k}\frac{\gamma_k}{\gamma_{k+1}}(\vy_k - \vy_{k-1}) - \frac{\Gamma_{k-1}}{\Gamma_k}\vy_k + \frac{\gamma_k}{\gamma_{k+1}}\frac{\mu}{4L}\vx_k$
    \ENDFOR
  \end{algorithmic}
\end{algorithm}

We consider \cref{alg:fast_gradient_method}, with a sequence $\{\gamma_k\}$ that satisfies $\gamma_0 = 1$ and  $2L+\Gamma_k\mu/2 = 2L\gamma_{k+1}^2/\Gamma_{k+1}$ with $\Gamma_k:= \sum_{i=0}^k \gamma_k$. This algorithms is an instantiation of the fast gradient method from \citep{devolder2014first}, in the case of unconstrained optimization. Originally formulated with three sequences and two projections, we simplify the method into two recursively defined sequences in \cref{app:sec:algorithm_simplification}.

As required by \cref{prop:link_beween_inexact_oracles}, we assume \cref{alg:fast_gradient_method} uses $(\zeta^2,\alpha=0)$-inexact gradient oracles, i.e. only an \textit{absolute} error is present. In a federated setting, this corresponds to $(G,B)$-Heterogeneity assumptions with $B=0$, common hypothesis in the literature (e.g. \citet{karimireddy2020byzantine}).

\begin{theorem}[Fast Gradient Method w. Inexact Oracles]
\label{thm:NAG_Inexact_Oracles}
    \cref{alg:fast_gradient_method} applied to a $\mu$-strongly convex and $L$-smooth function $\cL_{\cH}$, with a $(\zeta^2,\alpha)$-inexact gradient oracle, generates a sequence $(\vy_k)_{k\ge 1}$ satisfying, for $\alpha = 0$,
    \begin{align*}
    \cL_{\cH}(\vy_k) - \cL_{\cH}^* \le &\min\left(\frac{8LR}{k^2}, 2LR\exp\left(-\frac{k}{4}\sqrt{ \frac{L}{\mu} }\right)\right)\\ 
    &+ \left(1 + \sqrt{\frac{L}{\mu}}\right)\frac{3\zeta^2}{2\mu},
    \end{align*}
    where $R := \frac{1}{2}\|\vx_0 - \vx^*\|^2$.
\end{theorem}
\textit{Proof:} This follows from combining Theorem 7 in \citet{devolder2014first} with \cref{prop:link_beween_inexact_oracles}. 

\begin{corollary}
\label{cor:acceleration}
    Assume $(G,B=0)$-Heterogeneity (\cref{asmpt:GB_heterogeneity}), that $\cL_{\cH}$ is $\mu$-strongly convex and $L$-smooth, and that the robust aggregation rule $F$  in \cref{alg:reduction} is  $(f,\nu=\cO(\nicefrac{f}{n-2f}))$-robust. Then, \cref{alg:fast_gradient_method} using inexact oracle from \cref{alg:reduction}, requires
    $$
    k = \cO(\sqrt{\kappa}\log(\nicefrac{1}{\eps}))
    $$
    iterations to ensure
    \[
    \cL_{\cH}(\vy_k) - \cL_{\cH}(\vx^*) \le \cO\left( \eps + \frac{\sqrt{\kappa}}{\mu}\frac{f}{n-2f}G^2 \right).
    \]
\end{corollary}
The above result shows that \cref{alg:fast_gradient_method} enjoys an \textit{accelerated rate} towards the asymptotic error, i.e. the iteration complexity to reach a given error scales as the square root of the condition number $\kappa$, rather than linearly. To our knowledge, this is the first proof of such acceleration in the Byzantine setting.

\paragraph{Tradeoff: acceleration versus asymptotic error.}  The convergence radius in \cref{cor:acceleration} is suboptimal, inflated by $\sqrt{\kappa}$ factor w.r.t.~the lower bound (\ref{eq:lb_functionvalue}). This significant factor already appear in the acceleration results using $(\delta,L,\mu)$-oracles, and are unavoidable using such oracles. Indeed, \citet{devolder2014first} show that, since gradients of non-smooth functions are $(\delta,L,\mu)$-oracles, algorithms cannot converge arbitrarily fast under this assumption. Leveraging lower bounds on non-smooth Lipschitz functions, they show that any accelerated convergence rate under $(\delta,L,\mu)$ oracle must be traded off against a larger asymptotic error. Gradient Descent converges to the optimal neighborhood, but any faster rate suffers from a larger asymptotic error. 

Their lower bound does does not directly extend to our setting, since we consider smooth losses. This raises the question of whether the unavoidable tradeoff between acceleration and asymptotic error also holds for $(\zeta^2,\alpha)$-inexact gradient oracles or in the Byzantine setting, which we leave for future work.





\section{Gradient Descent under Similarity}
\label{sec:similarity}
We now propose a non-accelerated scheme that still reduces the number of communication rounds in the distributed algorithm. To do so, we leverage the so-called \textit{optimization under similarity} approach \citep{shamir2014communication, hendrikx2020statistically, kovalev2022optimal}, which assumes a slightly different setup.

\subsection{Setting}

We assume the server has unrestricted access to a proxy $\hat{\cL}$ of the global loss function $\cL_{\cH}$. This proxy may come from the empirical risk computed on a small toy dataset, or from a trusted client that using its own loss function as the proxy, i.e. $\hat{\cL} = \cL_1$. In this case, the server can leverage this cheap access to $\hat{\cL}$ (compared to $\cL_{\cH}$, which requires global communications) to improve the conditioning of the optimization problem. More specifically, the server solves the following problem at each iteration:
\begin{align}
    \label{eq:prox_update}
    \vx_{k+1} \approx \argmin_{\vx \in \R^d} \Big\{ &\overbrace{\hat{\cL}(\vx)}^{\text{\tiny Proxy Loss}} +  \langle \nabla\overbrace{(\cL_{\cH} - \hat{\cL})}^{\text{\tiny small Hessian}}(\vx_k), \vx \rangle \notag\\ 
    &+ \frac{1}{2\eta}\|\vx - \vx_k\|^2\Big\}.
\end{align}
This update can be seen from two perspectives. The mirror descent view \citep{shamir2014communication,hendrikx2020statistically} treats it as a mirror descent step on $\cL_{\cH}$ with step size $1$, and mirror map $\hat{\cL} + \frac{1}{2\eta}\| \cdot\|^2$. It can be shown that this problem satisfies relative smoothness assumptions~\citep{bauschke2017descent}, and thus enjoys linear convergence guarantees. Alternatively, the prox gradient perspective~\citep{kovalev2020optimal, woodworth2023two} rewrites the objective as $\cL_{\cH} = \hat{\cL} + \cL_{\cH} - \hat{\cL}$. %
One can then use a proximal gradient algorithm (a.k.a. \emph{forward-backward}), by taking a backward step on $\hat{\cL}$ (thus not depending on its smoothness), and a forward step on $\hat{\cL} - \cL_{\cH}$ (thus benefiting from its small smoothness constant if Hessians are close enough). Strong convexity just needs to hold for the global objective, also leading to a linear convergence rate. 

These manipulations modify the geometry of the problem, which is especially efficient when loss functions are ``similar" enough, namely when the Hessian of $\hat{\cL}$ is close to that of $\cL_{\cH}$. In both cases, they allow using significantly larger step sizes (from~$\frac{1}{L}$ to $\frac{1}{\Delta}$), hence significantly speeding up the optimization.

\subsection{Formal Assumptions}

We now state a similarity assumption between the proxy loss and the global loss. We also introduce a similarity assumption among clients' losses, which can be used to derive the similarity assumption between the proxy and the global loss, the $(G,B)$-heterogeneity condition, and a tighter characterization of the gradient inexactness.

\begin{assumption}[Hessian Similarity]
\label{asmpt:Hessian_Similarity}
    The global loss $\cL_{\cH}$ and the proxy loss $\hat{\cL}$ are said to have $\Delta$-similar Hessians if they are twice differentiable and
    \[
    \forall \vx \in \R^d, \|\nabla^2 \cL_{\cH}(\vx) - \nabla^2 \hat{\cL}(\vx)\|_{op} \le \Delta.
    \]
\end{assumption}

\cref{asmpt:Hessian_Similarity} arises, for instance, when the proxy used is a client's loss, $\cL_{1} = \hat{\cL}$, and the clients satisfy the pairwise Hessian Similarity assumption.

\begin{assumption}[Pairwise Hessian Similarity]
\label{asmpt:Pariwise_Hessian_Similarity}
    All local honest loss functions $(\cL_i)_{i \in \cH}$ are twice differentiable, and satisfy 
    \[
    \forall i,j \in\cH,\forall \vx \in \R^d, \quad \|\nabla^2 \cL_{i}(\vx) - \nabla^2 \cL_{j}(\vx)\|_{op} \le \Delta.
    \]
\end{assumption}

Interestingly, this latter similarity assumption directly implies $(G,B)$-Heterogeneity:

\begin{proposition}
    \label{prop:similarity_to_GB_heter}
    Let $(\cL_i)_{i\in\cH}$ be convex functions. Then, \cref{asmpt:Pariwise_Hessian_Similarity} ensures that
    \begin{align*}
    \frac{1}{|\cH|}\sum_{i\in\cH}\|&\nabla \cL_i(\vx)-\nabla \cL_{\cH}(\vx)\|^2 \le\\ 
    &\frac{2}{|\cH|}\sum_{i\in\cH}\|\nabla \cL_i(\vx^*)\|^2 + 4\Delta \langle \nabla \cL_{\cH}(\vx), \vx - \vx^*\rangle.
    \end{align*}
    Which, if additionally $\cL_{\cH}$ is $\mu$-strongly convex, implies $(G,B)$-Heterogeneity with 
    \[
    G^2 = \frac{2}{|\cH|}\sum_{i\in\cH}\|\nabla \cL_i(\vx^*)\|^2 \text{ and } B^2 = \frac{4\Delta}{\mu}.
    \]
\end{proposition}
We refer the reader to \cref{app:proof:similarity_to_GB_heter} for the proof.

\cref{prop:similarity_to_GB_heter} can be used to derive an alternative characterization of the gradient inexactness, on which we build our convergence guarantees.

\begin{corollary}
    Under \cref{asmpt:Pariwise_Hessian_Similarity}, \cref{alg:reduction}\ with a $(f,\nu)$-robust averaging rule (\cref{def:f_kappa_robustness}) yields:
    \begin{equation}
    \label{eq:inexact_gradient_inner_product}
    \!\!\!\|\tilde \nabla \cL_{\cH}(\vx) - \nabla \cL_{\cH}(\vx)\|^2 \le \zeta^2 + \alpha \mu \langle \nabla \cL_{\cH}(\vx), \vx - \vx^*\rangle,
    \end{equation}
    where $\alpha = \nu\frac{4\Delta}{\mu}$ and $\zeta^2 = \nu\frac{2}{|\cH|}\sum_{i \in \cH}\|\nabla \cL_i(\vx^*)\|^2$.
\end{corollary}

    We believe the characterization of oracle inexactness in  \cref{eq:inexact_gradient_inner_product} is novel and of independent interest, as it captures the loss function's growth directions more precisely than \cref{def:inexact_oracle}. Using as alternative parametrization $\tilde{\alpha} \leftarrow \alpha \mu = 4\Delta$ could enable the study of inexact gradient methods without requiring the strong convexity.

\subsection{Convergence Results}
 We now state the complete Byzantine-robust distributed algorithm along with its convergence guarantees.
 
\begin{algorithm}[h]
  \caption{Preconditioned Inexact Gradient w. Similarity}
  \label{alg:FedProxyProx}
  \begin{algorithmic}
    \STATE {\bfseries Input:} $F$, $\eta$, $\vx_0$, $\hat{\cL}$
     \FOR{$k=1$ {\bfseries to} $K$}
     \STATE Sample $\tilde\nabla \cL_{\cH}(\vx_k)$ using \cref{alg:reduction}.
    \STATE Compute $\vx_{k+1} \approx \argmin_{\vx \in \R^d} \phi_{k}(\vx)$\\
    \STATE where $\phi_k(\vx) = \hat{\cL}(\vx) +  \langle \tilde{\nabla} \cL_{\cH}(\vx_k) - \nabla \hat{\cL}(\vx_k), \vx \rangle + \frac{1}{2\eta}\|\vx - \vx_k\|^2$.
    \ENDFOR
  \end{algorithmic}
\end{algorithm}


\begin{theorem}
\label{thm:proxyprox_conv}
    Assume $\cL_{\cH}$ is $\mu$-strongly convex and $L$-smooth. Consider \cref{alg:FedProxyProx}, where the proxy loss $\hat{\cL}$ and the global loss $\cL_{\cH}$ satisfy \cref{asmpt:Hessian_Similarity}, and the inexact oracles $\tilde{\nabla}\cL_{\cH}(\vx_k)$ satisfy \cref{eq:inexact_gradient_inner_product} with $\alpha < \nicefrac{1}{8}$.

    Suppose $\eta \le (\Delta + 8c/\mu)^{-1}$, and the proximal step is computed with precision
    \begin{equation}
    \label{eq:precision_prox}
    \|\nabla \phi_k(\vx_{k+1})\|^2 \le c\|\vx_{k+1} - \vx_k\|^2 + E^2.
    \end{equation}
    Define $\beta_k = (1 + \frac{\eta \mu}{8})^{k}$, and $B_K = \sum_{k=0}^K \beta_k$. Then, for the weighted average $\hat{\vx}_k := \frac{1}{B_K}\sum_{k=0}^K \beta_k \vx_k$, we have
    \begin{align*}
    \cL_{\cH}(\hat\vx_{K}) - \cL_{\cH}(\vx^*) \le & \left(1 + \frac{\eta \mu}{8}\right)^{1-K}\!2\!\left(\!\frac{1}{\eta} + L\!\right)\!\|\vx_0 - \vx^*\|^2 \\
    & + 16\frac{\zeta^2 + E^2}{\mu}.
    \end{align*}
\end{theorem}
Importantly, in \cref{alg:FedProxyProx}, each $\vx_{k+1}$ is obtained by solving an optimization subroutine on the loss $\phi_k$ which the server can easily access. \cref{eq:precision_prox} captures the error of this subroutine, and both $c$ and $E^2$ can be made arbitrarily small by allocating more computational budget on the server side, i.e., without involving more communication or requiring more homogeneity or honest nodes. Notably, both $c$ and $E^2$ can be easily checked in practice. This subproblem can also be warm-started, starting the minimization from $\vx_k$, significantly reducing the computational burden.

The proof of \cref{thm:proxyprox_conv} extends the one of \citet{woodworth2023two}, who analyzed gradients method under similarity with only approximate proximal steps. In \cref{app:proof:proxyprox}, we introduce a new Lyapunov function to handle $(\zeta^2,\alpha)$-inexact gradients. We now instantiate \cref{thm:proxyprox_conv} in the Byzantine setting.

\begin{corollary}
    Let $(\cL_i)_{i \in \cH}$ satisfy \cref{asmpt:Pariwise_Hessian_Similarity}, and let $(G^2,B^2)$ be as in \cref{prop:similarity_to_GB_heter}. Assume $\cL_{\cH}$ is smooth and $\mu$-strongly convex. Let \cref{alg:FedProxyProx} use \cref{alg:reduction} with an $(f,\nu=\cO(\nicefrac{f}{n-2f}))$ robust-aggregation rule (\cref{def:f_kappa_robustness}) to compute $\tilde{\nabla}\cL$. Assume also that $\frac{f}{n-2f}\frac{\Delta}{\mu} \le C$, where $C$ is a numerical constant.
    
    Then, for a suitable step size $\eta = \nicefrac{1}{2\Delta}$, and a sufficiently small proximal error (\ref{eq:precision_prox}), \cref{alg:FedProxyProx} requires  
    \[\mathcal{O}\left(\frac{\mu}{\Delta}\log(1/\eps)\right)\] 
    communication rounds to guarantee
    \[
    \cL_{\cH}(\hat{\vx}_K) - \cL_{\cH}^* \in \mathcal{O}\left(\frac{f}{n-2f}\frac{G^2}{\mu} + \eps\right).
    \] 
\end{corollary}

\begin{figure}[b]
\centering
\includegraphics[width=\columnwidth]{mnist_logreg_mnist_n_21_f_1_dirichlet_niid_5.0_wd_0.01_TrMean_NNM_comparison_train_loss.pdf}
\caption{Distributed Logistic Regression, with $20$ honest and $1$ Byzantine client in the mildly heterogeneous setting ($\beta=5$), with a $l2$-regularization $\mu = 10^{-2}$. Train loss $\cL_{\cH}(\vx_k) - \cL_{\cH}^*$ under A Little Is Enough (ALIE) attack.}
\label{plot:logistic_heterogeneous_loss}
\end{figure}


\begin{remark}
     Note that $\nicefrac{\Delta}{\mu}$ can be much smaller than $\kappa = \nicefrac{L}{\mu}$, the original condition number of the global loss $\cL_{\cH}$, thus significantly speeding-up optimization. Indeed, assume that $\cL_i(\vx) := \frac{1}{|\cD_i|}\sum_{a \in \cD_i}l(a,\vx) + \lambda/2\|\vx\|^2$, where each dataset $\cD_i$ consist of $m$ i.i.d. samples from the same distribution $\cD$. Assume $\|\nabla^2_{\vx}l(a,\vx)\|_{op} \le L$. %
     Then w.l.o.g., $\|\nabla^2 \cL_{\cH}(\vx)\|_{op} \le L$, while Hoeffding's inequality for matrices~\citep{tropp2015introduction} yields, with probability at least $1-\delta$, 
    \[
    \Delta = \|\nabla^2 \cL_1(\vx) - \nabla^2 \cL_{\cH}(\vx)\|_{op} \le \mathcal{O}\left(\sqrt{\frac{L^2\log(d/\delta)}{m}}\right).
    \]
    It follows $\frac{\mu}{\Delta} \le \mathcal{O}\left(\frac{\kappa}{\sqrt{m}}\log(d/\delta)\right)$. Thus, the iteration complexity of \cref{alg:FedProxyProx} to reach a given precision can be reduced by up to $\mathcal{O}(\sqrt{m})$ compared to Byzantine-robust D-GD, ignoring logarithmic factors. When local datasets are large, this yields a huge improvement in the communication complexity. For refined and stronger control of $\Delta$, including better dependencies on $m$, see \citet{hendrikx2020statistically}.
\end{remark}



\section{Experiments}
\label{sec:experiments}
\begin{figure*}[t]
\centering
\includegraphics[width=0.9\textwidth]{mnist_logreg_mnist_n_21_f_1_dirichlet_niid_1.0_wd_0.001_TrMean_NNM_comparison_with_ref_train_loss.pdf}
\caption{Impact of the learning rate $\eta$ on PIGS, in a highly heterogeneous setting $(\beta=1)$ with $20$ honest clients and either with $1$ Byzantine performing IPM or ALIE attack, or without Byzantine client. $l2$-regularization with $\mu = 10^{-3}$.}
\label{plot:logistic_highly_heterogeneous_loss}
\end{figure*}
\subsection{Setup}
\textbf{Datasets and models.} We consider a Logistic Regression task on the MNIST dataset \citep{lecun1998gradient}. To simulate heterogeneity among clients, we follow the Dirichlet partitioning method from \citet{hsu2019measuring, allouah2023fixing}: each client’s dataset class proportions are sampled from a Dirichlet distribution with parameter $\beta$. Lower $\beta$ corresponds to higher heterogeneity. We refer to $\beta=1$ as \textit{highly heterogeneous} (\cref{plot:logistic_highly_heterogeneous_loss}) and to $\beta = 5$ (\cref{plot:logistic_heterogeneous_loss}) as \textit{mildly heterogeneous}. We also provide experiments with a random split of the training data among clients, referred to as the \textit{i.i.d.} setting (\cref{plot:logistic_iid_loss}). We reefer to \cref{app:experiments} for the test accuracy curves.

\textbf{Optimization Algorithms.} We compare Distributed Nesterov's Accelerated Gradient (D-NAG), Distributed Gradient Descent (D-GD) and PIGS. In all methods, inexact gradient are computed using \cref{alg:reduction}, with gradients aggregated using Nearest Neighbors Mixing (NNM) combined with Coordinate-Wise Trimmed Mean (CWTM). The proximal step in PIGS is solved approximately using L-BFGS.

\textbf{Attacks and parameters.} We tune optimizer parameters via hyperparameter search. The attacks implemented are variants of A Little Is Enough (ALIE) \citep{baruch2019little} and Inner Product Manipulation \citep{xie2020fall} where the attack is scaled using a line search to maximize impact. Our code, accessible \href{https://github.com/renaudgaucher/Inexact-gradient-acceleration-and-similarity}{here}, is built on top of the ByzFL library \citep{gonzalez2025byzflresearchframeworkrobust}.

\subsection{Analysis}

\begin{figure*}[t]
\centering
\includegraphics[width=0.9\textwidth]{mnist_logreg_mnist_n_25_f_5_iid_None_wd_0.001_TrMean_NNM_comparison_with_ref_train_loss.pdf}
\caption{Impact of the learning rate $\eta$ on PIGS,i.i.d. setting with $20$ honest clients and either with $5$ Byzantine performing IPM or ALIE attack, or without Byzantine client. $l2$-regularization with $\mu = 10^{-3}$.}
\label{plot:logistic_iid_loss}
\end{figure*}

\textbf{Communication Complexity.} In all experiments, PIGS achieves significantly better communication complexity than D-NAG and D-GD. For instance, in \cref{plot:logistic_heterogeneous_loss}, PIGS requires only $5$ iterations to go below the asymptotic error, while Nesterov's acceleration requires $41$ iterations, and D-GD $135$ iterations. As predicted by our theory, the faster rate of PIGS is derived from the stable use of larger step sizes when the proxy loss is similar to the global loss. For instance, in the highly heterogeneous setting (\cref{plot:logistic_highly_heterogeneous_loss}) $\eta=2$ seems the highest possible step size when $f=0$, while in the i.i.d. setting (\cref{plot:logistic_iid_loss}) $\eta = 8.$ still ensures convergence.

\textbf{Stability.} As shown in \cref{plot:logistic_highly_heterogeneous_loss,plot:logistic_iid_loss}, neither ALIE nor IPM attacks destabilize PIGS's convergence when the learning rate $\eta$ is large: it is more stable under the tested Byzantine attacks than without adversaries. This is obviously a limitation of current attacks, which focus on \textit{slowing} or \textit{biasing} the optimization, rather than destabilizing it. 

\textbf{Asymptotic Error.} As shown in \cref{plot:logistic_heterogeneous_loss} the three optimization algorithms converge to the same asymptotic loss. This raises the question of whether the $\sqrt{\kappa}$ multiplicative factor in the asymptotic error of \cref{thm:NAG_Inexact_Oracles} can be alleviated our not. Since existing attacks have been designed mostly to tackle Gradient Descent, this might be an artifact of their design. In order to make the $\sqrt{\kappa}$ multiplicative factor appear on the sub-optimality, attack should stress the weak points of accelerated methods, namely their instabilities. We leave this for future work.




\section{Additional Discussion on the Tightness of the Abstraction}




Designing Byzantine-robust algorithms via \cref{alg:reduction} and \cref{lemma:Byz_to_Inexact_Gradients} abstracts away the distributed nature of the problem. This raises the question of whether this abstraction is always tight, and whether it recovers the guarantees obtainable in the full distributed setting.

\textbf{Deterministic Case.} As pointed out in \cref{subsec:tightness}, the achievable error and the breakdown requirement of the abstraction are optimal, in the sense that GD achieves both an optimal asymptotic error and has an optimal breakdown requirement. On the other hand, the optimal convergence rate in the Byzantine context is unknown, and so is the optimal convergence rate with $(\zeta^2,\alpha)$-inexact oracles, even if several works have studied the impact of $\alpha$ on the convergence rate of Gradient methods under $\zeta^2=0$ \citep{de2020worst, gannot2022frequency}. 

\textbf{Stochastic Case.} It has been known since \citet{karimireddy2021learning} that using abstractions such as \cref{alg:reduction} cannot ensure reaching the optimal asymptotic error when the clients send unbiased stochastic oracles of their local loss gradient. Indeed, \citet{karimireddy2021learning} demonstrate that memory less optimization algorithms, which do not track stochastic oracles as deriving from specific clients, can only reach an optimization error in $\Omega(\frac{f}{n-f}\frac{\sigma^2}{\mu})$, where $\sigma^2$ is the scale of the stochastic noise. \\

\section{Conclusion}
In this work, we proposed to abstract away the problem of Byzantine Optimization by formalizing it as an inexact gradient oracle problem. Our approach opens up a new systematic approach to study and design Byzantine resilient algorithms. We leverage this to propose the first accelerated scheme in this context, as well as an optimization under similarity method. Those two methods have better communication complexity than previous ones, as evidenced experimentally and theoretically.








\section*{Impact Statement}

This paper presents work whose goal is to advance the field of Machine
Learning. There are many potential societal consequences of our work, none
of which we feel must be specifically highlighted here.

\bibliography{biblio}
\bibliographystyle{icml2026}


\newpage
\appendix
\onecolumn






\section*{Appendices}
The appendices are organized as follows:
\begin{enumerate}
    \item In \cref{app:sec:robust_aggregation}, we recall the robustness coefficients of standard aggregation rules. We also extend the mixing strategy of \citet{allouah2023fixing} using the F-Robust Gossip framework of \citet{gaucher2025unified}, enabling the design of tighter aggregation rules.
    \item In \cref{app:proofs}, we present the proofs of our theoretical results:
    \begin{enumerate}
        \item \cref{app:proofs_standard_tools} reviews the standard optimization tools used throughout the proofs.
        \item \cref{app:proof:proxyprox} provides the proof of \cref{thm:proxyprox_conv}.
        \item \cref{app:sec:algorithm_simplification} shows how the fast gradient method of \citet{devolder2014first} can be reformulated as a two-sequence algorithm.
    \end{enumerate}
    \item In \cref{app:experiments}, we present additional experiments, including the test accuracy curves corresponding to the main results.
\end{enumerate}

\section{Robust Aggregation Rules}
\label{app:sec:robust_aggregation}
We recall the list of $(f,\nu)$-robust aggregation rules established by \citet{allouah2023fixing}. 

\begin{table}[h]
    \centering
    \renewcommand{\arraystretch}{1.2}
    \begin{tabular}{c c c c c | c}
    \toprule
         Aggregation &  CWTM & Krum & GM & CWM & Lower Bound\\ %
    \midrule
         $\kappa$ & $\frac{6f}{n-2f}\left(1 + \frac{6f}{n-2f}\right)$ & $6\left(1 + \frac{6f}{n-2f}\right)$ & $4\left(1 + \frac{f}{n-2f}\right)^2$ & $4\left(1 + \frac{f}{n-2f}\right)^2$ & $\frac{f}{n-2f}$\\ %
    \bottomrule
         
    \end{tabular}
    \caption{Robustness coefficients of the coordinate-wise trimmed mean (CWTM), Krum, the geometric median (GM) and the coordinate-wise median (CWM)}
\end{table}

\citet{allouah2023fixing} suggest to improve aggregation rules by pre-processing the input values through a pre-aggregation step, and specifically the Nearest Neighbors Mixing (NNM) strategy. Which we extend in a formal definition:
\begin{definition}
    A map $F_1:\R^{n\times d} \rightarrow \R^{n\times d}$ is said to be a $(f, \delta)$ robust mixing if, when $|\cB| \le f$, and given $F_2$ a $(f,\nu)$-robust aggregation rule, $F_2 \circ F_1$ is a $(f, \tilde \nu)$-robust aggregation rule with
    \[
    \tilde \nu = \delta (1 + \nu).
    \]
\end{definition}
Below are listed few robust mixing, which are can be build generally using the robust-gossip framework, as we show in \cref{app:sec:robst_gossip_to_robust_mixing}.

\begin{table}[h]
    \centering
    \renewcommand{\arraystretch}{1.2}
    \begin{tabular}{c c c c | c}
        \toprule
         Pre-Aggregation & CS$_+$-RG & NNM/NNA & F-RG  & Lower Bound\\
         \midrule
         $\delta$ &   $\frac{4f}{n-f}$ & $\frac{8f}{n-f}$  & $\frac{2\rho f}{n-f}$&  \\
         Breakdown & $\frac{1}{5}$ & $\frac{1}{9}$ & $\frac{1}{2\rho + 1}$ &  $\frac{1}{3}$ \\
         \bottomrule
    \end{tabular}
    \caption{Robust Pre-Aggregation / Mixing Strategies. F in F-RG is a $(\rho, b=\frac{f}{n-f})$-robust summand. }
\end{table}



\subsection{From Byzantine-Robust Gossip to Byzantine-Robust Distributed Learning}
\label{app:sec:robst_gossip_to_robust_mixing}
\begin{algorithm}[h]
  \caption{Robust-Gossip based Pre-Aggregation}
  \label{app:alg:F_RG_preaggregation}
  \begin{algorithmic}
    \STATE {\bfseries Input:} $F_2$ $(\rho, \frac{f}{n-f})$-robust summand rule, $F_1$ $(f,\nu)$-robust aggregation rule, vectors input $(\vx_i)_{i\in [n]}$.
    \FOR{Client $i \in [n]$}
     \STATE Client Send $\vx_i$ to Server
    \ENDFOR
    \STATE Server Update $\vy_i = \vx_i - F( (\vx_i - \vx_j, \frac{1}{n-f})_{j\in [n]})$
  \end{algorithmic}
\end{algorithm}


Nearest Neighbor Mixing (NNM) can be seen as an emulation of the Nearest Neighbor Averaging (NNA) algorithm \citep{farhadkhani2023robust}. As such, the mixing strategy constitutes an intrinsically decentralized aggregation scheme, which can be further generalized within the Robust Gossip framework \citep{gaucher2025unified}. 
More precisely, if the server simulates a single iteration of a $F$-Robust Gossip \citep{gaucher2025unified}, with $F$ a $(b,\rho)$-robust summand, the resulting update satisfies the mixing property required by NNM. This allows to design even tighter pre-aggregation strategy. 


$F$-Robust Gossip (or $F$-RG) is a decentralized algorithm, in the sense that all clients are seen as vertices in a communication graph $\cG$, in which they communicate synchronously with their neighbors through the edges. The algorithm relies on so-called $(b,\rho)$-robust summand functions, defined as follows:

\begin{definition}[$(b,\rho)$--robust summation] Let $b, \rho \ge 0$. An aggregation rule $F:(\R_+ \times \R^d)^n \rightarrow \R^d$ is a \emph{$(b, \rho)$--robust summation} if, for any vectors $(\vz_i)_{i \in [n]}\in (\R^d)^n$, any weights $(\omega_i)_{i \in [n]} \in \R_+^n$ and any $S \subset [n]$ such that $\sum_{i \in \overline{S}} \omega_i \le b$ (where $\overline{S} := [n]\backslash S$),
\[
\bigg\|F\big((\omega_i, \vz_i)_{i \in [n]}\big) - \sum_{i\in S} \omega_i \vz_i\bigg\|^2 \le \rho b \sum_{i\in S} \omega_i \|\vz_i\|^2.
\]
\end{definition}

Given $F$ a $(b,\rho)$-robust summand, and weights $(w_{ij})_{i,j\in [n]}$, one step of $F$-RG updates the client's parameter $(\vx_i)_{i \in [n]}$ according to
\[
\forall i \in \cH,\quad  \vy_i = \vx_i - \eta F( (\vx_i - \vx_j, w_{ij})_{j \in [n]}).
\]

One particular instance of $F$-RG simulated by a central server in a server–client architecture is the Nearest Neighbor Mixing (NNM) pre-aggregation scheme \citep{allouah2023fixing}. This method corresponds alternatively to NNA \citep{farhadkhani2023robust} or to $F$-RG with weights $w_{ij} = 1/(n-f)$ and step size $\eta = 1$ when $F$ is chosen as the geometrically trimmed sum (GTS), formally:
\[
\mathrm{GTS}\left((\vz_1,\nicefrac{1}{n-f}),\ldots, (\vz_n, \nicefrac{1}{n-f}) \right) := \frac{1}{n-f}\sum_{n-f \text{ smallest } \|\vz_i\|}\vz_i.
\]



\begin{proposition}
    Let's $\cG$ be the fully-connected graph on $[n]$ with uniform weights on the edges $\frac{1}{n-f}$. Let $F_1$ be a $(f,\nu)$-robust aggregator, and $F_2$ be a $(\rho,b)$ robust summand, with $b=\frac{f}{n-f}$. Then $F_1 \circ(F_2\text{-}RG)$ as in \cref{app:alg:F_RG_preaggregation} is a $(\tilde{\kappa},f)$-robust aggregator, with 
    \[
    \tilde{\kappa} = \delta (1+\kappa)\quad \text{where} \quad \delta := \frac{2\rho f}{n-f}.
    \]
    We refer to $\delta$ as the mixing contraction factor.
\end{proposition}


\begin{proof}
    
    \citet{gaucher2025unified} show that, for graph $\cG$ of spectral gap $\gamma$, $\rm F_2-RG$, with $F_2$ a $(b,\rho)$-robust aggregator with step size $\eta = \mu_{\max}^{-1}(\cG_{\cH})$, ensures that, for $(\vy_i)_{i \in [n]} = \mathrm{F_2-RG}(\vx_i, i\in [n])$, and $b < \mu_2(\cG_{\cH})$
    \begin{equation}
    \label{app:frg_contraction}
    \frac{1}{|\cH|}\sum_{i\in \cH} \|\vy_i - \overline{\vx}_{\cH}\|^2 \le (1 - \gamma(1 - \delta))\frac{1}{|\cH|}\sum_{i \in \cH}\|\vx_i - \overline{\vx}_{\cH}\|^2,
    \end{equation}
    where $\delta = \frac{2\rho b}{\mu_2(\cG_{\cH})}$.
    
    Thus, plugging the $\vy_i$ in $F_2$ leads to 
    \begin{align*}
        \|F_2(\vy_1,\ldots,\vy_n) - \overline{\vx}_{\cH}\|^2 &\le (1 + \eps^{-1}) \|F_2(\vy_1,\ldots,\vy_n) - \overline{\vy}_{\cH}\|^2 + (1 + \eps) \|\overline{\vy}_{\cH} - \overline{\vx}_{\cH}\|^2 && \text{Young's inequality}\\
        &\le \kappa (1 + \eps^{-1}) \frac{1}{|\cH|}\sum_{i\in\cH}\|\vy_i - \overline{\vy}_{\cH}\|^2 + (1 + \eps) \|\overline{\vy}_{\cH} - \overline{\vx}_{\cH}\|^2 && (f,\kappa)\text{-robustness}\\
        &= (1+\kappa)\left(\frac{1}{|\cH|}\sum_{i\in\cH}\|\vy_i - \overline{\vy}_{\cH}\|^2 + \|\overline{\vy}_{\cH} - \overline{\vx}_{\cH}\|^2 \right) \quad \text{for } \eps = \kappa\\
        &\le (1+\kappa)(1 - \gamma(1-\delta)) \frac{1}{|\cH|}\sum_{i \in |\cH|}\|\vx_i - \overline{\vx}_{\cH}\|^2\ , 
    \end{align*}   
    where the last inequality stems from \cref{app:frg_contraction} and a bias-variance decomposition.

    Thus, if the central server emulates one step of $F_2$-RG on a fully connected graph, with uniform weights $\frac{1}{|\cH|}$, we have that $\gamma=1$, $\mu_{\max}(\cG_{\cH}) = \mu_{\min}(\cG_{\cH}) = \eta^{-1} = 1$, and thus $\delta = \frac{2\rho f}{n-f}$, and use $F$ on the outputs (which corresponds to NNA for G$=$Trimmed Sum), we have that
    \[
    \tilde{\kappa} = \delta(1+\kappa).
    \]
\end{proof}
\begin{remark}
    Considering sparser graphs with large spectral gap (e.g. expander) leads to similar results, while reducing the computational cost of $O(n^2 d)$, without loosing much on $\gamma$ and $\delta$.
\end{remark}












\section{Proofs}
\label{app:proofs}
\subsection{Introduction of Standard Optimization tools}
\label{app:proofs_standard_tools}
In the proof, we  use several standard optimization tools, that we recall hereafter.
\begin{definition}[Bregman Divergence]
    The Bregman divergence of a real-valued function $f$ is defined as 
    \[
    D_f(\vx,\vy) := f(\vx) - f(\vy) - \langle \nabla f(\vy), \vx - \vy \rangle.
    \]
\end{definition}
The Bregman divergence has the following properties:
\begin{proposition}
    \begin{enumerate}
        \item A function $f$ is $\mu$-strongly convex and $L$-smooth, if and only if
        \begin{equation}
        \label{app:eq:smoothness_strongcvx_bregman}
        \forall \vx,\vy \in \R^d, \qquad \frac{\mu}{2}\|\vx - \vy\|^2 \le D_f(\vx;\vy) \le \frac{L}{2}\|\vx-\vy\|^2.
        \end{equation}
        \item \textbf{Three points identity.} Let $\vx,\vy,\vz \in \R^d$, it holds:
        \begin{equation}
        \label{app:eq:three_points_identity}
        D_{f}(\vx;\vy) + D_{f}(\vy;\vz) - D_{f}(\vx;\vz) = \langle  \nabla f(\vy) - \nabla f(\vz), \vy - \vx\rangle.
        \end{equation}
    \end{enumerate}
\end{proposition}

Furthermore, we will  use the following property of strongly convex functions:

\begin{proposition}
    Let $f:\R^d \rightarrow \R$ be a $\mu$-strongly convex $L$-smooth loss function, then:
    \begin{equation}
        \label{app:eq:cocoaxivity_strong_convexity}
            \langle\nabla f(\vx) - \nabla f(\vy), \vx - \vy\rangle \le \frac{1}{\mu}\|\nabla f(\vx) - \nabla f(\vy)\|^2.
        \end{equation}
\end{proposition}
\begin{proof}
        We use that the Fenchel transform of $f$, defined as $f^*:\vu \rightarrow \sup_{\vx \in \R^d}\{\langle \vx,\vu\rangle - f(\vx)\}$, is $1/\mu$ smooth. 
        This ensures that 
        \[
        \forall \vv, \vu \in \R^d,\quad \langle \vu - \vv, \nabla f^*(\vu) - \nabla f^*(\vv)\rangle \le \frac{1}{\mu}\|\vu - \vv\|^2.
        \]
        Then, using Fenchel's identity, which here can be written as $\nabla f = (\nabla f^*)^{-1}$, and considering $\vx = \nabla f^*(\vu)$ and $\vy = \nabla f^*(\vv)$ yields the result.
\end{proof}

\subsection{Optimization under Similarity}
\label{app:proof:proxyprox}
We  show the following theorem, of which \cref{thm:proxyprox_conv} is a direct consequence.

\begin{theorem}
\label{app:thm:proxy_prox_inexact}
   Let $f:\R^d \rightarrow \R$ be a $\mu$-strongly convex and $L_f$-smooth loss function, written as $f = h + g$, where $h$ is $L_h$-smooth (i.e. $\|\nabla^2 h(\vx)\|_{op} \le L_h$). 
Assume that we have access to an inexact gradient oracle $\tilde \nabla h$ satisfying
\begin{equation}
\label{app:eq:inexact_gradient_scalar_product}
\forall \vx\in\R^d, \quad\|\tilde \nabla h(\vx) - \nabla h(\vx)\|^2 \le \zeta^2 + \alpha \mu \langle \nabla f(\vx), \vx - \vx^* \rangle,
\end{equation}
where \(\vx^* := \argmin_{\vx\in \R^d}f(\vx)\) and $\alpha < \nicefrac{1}{8}$.
   Consider $(\vx_{k})_{k\ge 0}$ such that for any $k$, given 
   \begin{equation}
   \label{app:eq:def_phi_k}
       \tilde{\phi}_k(\vx) := \langle \tilde{\nabla} h(\vx_k) + \nabla g(\vx_k), \vx \rangle + D_{g}(\vx;\vx_k) + \frac{1}{2\eta}\|\vx - \vx_k\|^2,
   \end{equation}
$\vx_{k+1}$ is chosen such that
\begin{equation}
\label{app:eq:error_in_prox}
\|\nabla \tilde{\phi} (\vx_{k+1})\|^2 \le c\|\vx_{k+1} - \vx_k\|^2 + E^2.
\end{equation}

Denote $\beta_k = (1 + \frac{\eta \mu}{8})^{k}$, and $B_K = \sum_{k=0}^K \beta_k$. Then for $\eta \le \frac{1}{L_h + 8c/\mu}$, and $\hat{\vx}_k := \frac{1}{B_K}\sum_{k=0}^K \beta_k \vx_k$, we have

\[
f(\hat\vx_{K}) - f(\vx^*) \le \left(1 + \frac{\eta \mu}{8}\right)^{1-K}2(\eta^{-1} + L_f)\|\vx_0 - \vx^*\|^2 + 16\frac{\zeta^2 + E^2}{\mu}.
\]
\end{theorem}

Note that we did lose some constant factors for the sake of clarity.

\subsubsection{Proof of \cref{app:thm:proxy_prox_inexact}}
This proof is an adaptation of the proof of \citet{woodworth2023two}, in which we modify the Lyapunov function to allow for the multiplicative error term (i.e $\alpha \neq 0$) in the inexact gradient assumption. 

We first show the following decomposition

\begin{lemma}
\label{app:lemma:proof_proxyprox_1}
Denote $\ve_k := \nabla h(\vx_k) - \tilde \nabla h(\vx_k)$, such that, by definition $\|\ve_k\|^2 \le \zeta^2 + \alpha \mu \langle \nabla f(\vx_k),\vx_k - \vx^*\rangle$. Then the following identity holds:
    \begin{align*}
        f(\vx_{k+1}) - f(\vx^*) &=  \left(\frac{1}{2\eta}\|\vx_{k}- \vx^*\|^2 - D_{h}(\vx^*;\vx_{k})\right) - \left(\frac{1}{2\eta}\|\vx_{k+1}- \vx^*\|^2 - D_{h}(\vx^*;\vx_{k+1}) \right) \\
        &- D_f(\vx^*;\vx_{k+1})\\
        & + \langle \nabla \tilde \phi_k(\vx_{k+1}) + \ve_k, \vx_{k+1} - \vx^* \rangle \\
        &+ D_h(\vx_{k+1};\vx_k) - \frac{1}{2\eta}\|\vx_{k}- \vx_{k+1}\|^2.
    \end{align*}
\end{lemma}
The first line is telescopic, the second line will be leveraged to ensure a strict decrease in the Lyapunov value, the third line corresponds to the error due to the approximate proximal step, and the erroneous gradients. Eventually, the last line corresponds to a discretization error.


\begin{proof}

    Computing $\nabla \tilde\phi_k(\vx_{k+1})$ using \cref{app:eq:def_phi_k} and $\nabla_\vx D_g(\vx,\vx_k) = \nabla g(\vx) - \nabla g(\vx_k)$, defining $\ve_k = \nabla h(\vx_{k}) - \tilde \nabla h(\vx_k)$ the error in the computation of $\nabla h(\vx_k)$, and rearranging terms yields
    \begin{align*}
        \nabla \tilde\phi_k(\vx_{k+1}) &= \tilde \nabla h(\vx_{k}) + \nabla g(\vx_{k}) + \nabla g(\vx_{k+1}) - \nabla g(\vx_k) + \frac{1}{\eta}(\vx_{k+1} - \vx_k)\\
        &= \underbrace{\tilde{\nabla} h(\vx_k) - \nabla h(\vx_k)}_{= - \ve_k} + \nabla h(\vx_k) - \nabla h(\vx_{k+1}) + \underbrace{\nabla h(\vx_{k+1}) + \nabla g(\vx_{k+1}) }_{= \nabla f(\vx_{k+1})} + \frac{1}{\eta}(\vx_{k+1}-\vx_k)
        \\
        \nabla f(\vx_{k+1}) &= \nabla h(\vx_{k+1}) - \nabla h(\vx_k)  + \nabla \tilde \phi_k(\vx_{k+1}) + \ve_k + \frac{1}{\eta}(\vx_{k}-\vx_{k+1})\ .
    \end{align*}

    Furthermore, the definition of the Bregman divergence implies 
    \[
    f(\vx_{k+1}) - f(\vx^*) = \langle \nabla f(\vx_{k+1}), \vx_{k+1}- \vx^*\rangle - D_f(\vx^*; \vx_{k+1}).
    \]
    Substituting the previous identity into the latter one yields
    \begin{align*}
        f(\vx_{k+1}) &- f(\vx^*)\\
        &= \langle \nabla f(\vx_{k+1}), \vx_{k+1}- \vx^*\rangle - D_f(\vx^*; \vx_{k+1})\\
        &= \left \langle \nabla h(\vx_{k+1}) - \nabla h(\vx_k)  + \nabla \tilde \phi_k(\vx_{k+1}) + \ve_k + \frac{1}{\eta}(\vx_{k}-\vx_{k+1})
        , \vx_{k+1}- \vx^*\right \rangle - D_f(\vx^*; \vx_{k+1})\\
        &=  D_{h}(\vx^*;\vx_{k+1}) + D_h(\vx_{k+1};\vx_k) - D_{h}(\vx^*;\vx_{k}) \\
        &+ \langle \nabla \tilde \phi_k(\vx_{k+1}) + \ve_k, \vx_{k+1} - \vx^* \rangle \\
        &- D_f(\vx^*;\vx_{k+1}) + \frac{1}{2\eta}\|\vx_{k}- \vx^*\|^2 - \frac{1}{2\eta}\|\vx_{k}- \vx_{k+1}\|^2 - \frac{1}{2\eta}\|\vx_{k+1}- \vx^*\|^2.
    \end{align*}

    Where we used the parallelogram identity,
    \[
    \frac{1}{\eta}\langle \vx_k - \vx_{k+1}, \vx_{k+1} - \vx^*\rangle = \frac{1}{2\eta}\|\vx_k  - \vx^*\|^2 - \frac{1}{2\eta}\|\vx_k - \vx_{k+1}\|^2 - \frac{1}{2\eta}\|\vx_{k+1} - \vx^*\|^2,
    \]
    and the three-point identity \cref{app:eq:three_points_identity},
    \[
    D_{h}(\vx^*;\vx_{k+1}) + D_h(\vx_{k+1};\vx_k) - D_{h}(\vx^*;\vx_{k}) = \langle  \nabla h(\vx_{k+1}) - \nabla h(\vx_k), \vx_{k+1} - \vx^*\rangle.
    \]
\end{proof}

We now control the error due to inexact oracles and approximate proximal step as follows:
\begin{lemma}
    \label{app:lemma:proof_proxyprox_2}
    For any $\eps >1$, the following inequality holds, for any $k\ge 0$,
    \[
    \langle \nabla \tilde \phi_k(\vx_{k+1}) + \ve_k, \vx_{k+1} - \vx^*\rangle \le \frac{c}{2\eps}\|\vx_{k+1} - \vx_k\|^2  + \frac{\alpha \mu }{2\eps} \langle\nabla f(\vx_k), \vx_k - \vx^* \rangle+ \eps\|\vx_{k+1}-\vx^*\|^2 + \frac{E^2 + \zeta^2}{2\eps}.
    \]
\end{lemma}


\begin{proof}
    Writing the inexact gradient assumption (\cref{app:eq:inexact_gradient_scalar_product}) and Young's inequality, we have for any $\eps >0$,
    \begin{equation}
    \label{app:eq:control_Young_inexact_gd}
        \langle \ve_k, \vx_{k+1} - \vx^*\rangle \le \frac{1}{2\eps}\left(\zeta^2 + \alpha \mu  \langle\nabla f(\vx_k), \vx_k - \vx^* \rangle\right) + \frac{\eps}{2}\|\vx_{k+1}-\vx^*\|^2.
    \end{equation}
    Using the assumption \cref{app:eq:error_in_prox} on the proximal step, namely
    \[
    \|\nabla \tilde \phi_k(\vx_{k+1}\|^2 \le c \|\vx_{k+1} - \vx_k\|^2 + E^2,
    \]
    and Young's inequality, we have that 
    \begin{equation}
    \label{app:eq:control_Young_prox}
        \langle \nabla \tilde \phi_k(\vx_{k+1}), \vx_{k+1} - \vx^*\rangle \le \frac{1}{2\eps}\left(c\|\vx_{k+1} - \vx_k\|^2 + E^2\right) + \frac{\eps}{2}\|\vx_{k+1}-\vx^*\|^2.
    \end{equation}
    Putting \cref{app:eq:control_Young_inexact_gd} and \cref{app:eq:control_Young_prox} together gives \cref{app:lemma:proof_proxyprox_2}.
\end{proof}

We now define a Lyapunov energy function, and show its decreasing.

\begin{lemma}
\label{app:lemma:proof_proxyprox_3}
    Let's consider the Lyapunov 
    \begin{equation}
    \label{app:eq:def_Lyapunov}
    \Gamma_{k} := \frac{1}{2\eta}\|\vx_k - \vx^*\|^2 - D_{h}(\vx^*;\vx_k) + \frac{\alpha \mu}{2 \eps} \langle \nabla f(\vx_k),\vx_k - \vx^*\rangle.
    \end{equation}
    Then, under the choice $\eps = \frac{\mu}{8}(1 + \sqrt{1 - 8\alpha})$, we have
    \begin{align*}
        \frac{1}{4}(f(\vx_{k+1}) - f(\vx^*)) 
        \le \Gamma_k - \left(1 +  \frac{\eta \mu}{8}\right)\Gamma_{k+1} + \frac{4\zeta^2}{\mu} + \frac{4E^2}{\mu}.
    \end{align*}
\end{lemma}
\begin{proof}
    Combining \cref{app:lemma:proof_proxyprox_1} and \cref{app:lemma:proof_proxyprox_2}, gives 
    \begin{align*}
        f(\vx_{k+1}) - f(\vx^*) &\le  \left(\frac{1}{2\eta}\|\vx_{k}- \vx^*\|^2 - D_{h}(\vx^*;\vx_{k})\right) - \left(\frac{1}{2\eta}\|\vx_{k+1}- \vx^*\|^2 - D_{h}(\vx^*;\vx_{k+1}) \right) \\
        &- D_f(\vx^*;\vx_{k+1})\\
        &+ \frac{c}{2\eps}\|\vx_{k+1} - \vx_k\|^2  + \frac{\alpha \mu }{2\eps} \langle\nabla f(\vx_k), \vx_k - \vx^* \rangle+ \eps\|\vx_{k+1}-\vx^*\|^2 + \frac{E^2 + \zeta^2}{2\eps}\\
        &+ D_h(\vx_{k+1};\vx_k) - \frac{1}{2\eta}\|\vx_{k}- \vx_{k+1}\|^2.
    \end{align*}
    
    Thus, adding on both sides $\frac{\alpha \mu}{2\eps}D_f(\vx^*;\vx_{k+1}) = \frac{\alpha \mu}{2\eps}[f(\vx^*) - f(\vx_{k+1}) + \langle \nabla f(\vx_{k+1}), \vx_k - \vx^* \rangle]$ and rearranging yields
    \begin{align*}
         \left(1 - \frac{\alpha \mu}{2 \eps}\right)\left(f(\vx_{k+1}) - f(\vx*)\right) &\le
         \left(\frac{1}{2\eta}\|\vx_{k}- \vx^*\|^2 - D_{h}(\vx^*;\vx_{k}) + \frac{\alpha \mu}{2 \eps} \langle \nabla f(\vx_k),\vx_k - \vx^*\rangle\right) \\
         &- \left(\frac{1}{2\eta}\|\vx_{k+1}- \vx^*\|^2 - D_{h}(\vx^*;\vx_{k+1}) + \frac{\alpha \mu}{2 \eps} \langle \nabla f(\vx_{k+1}),\vx_{k+1} - \vx^*\rangle \right)\\
         &- \left(1 - \frac{\alpha \mu}{2 \eps}\right) D_f(\vx^*;\vx_{k+1})\\
        & + \frac{c}{2\eps}\|\vx_{k+1} - \vx_k\|^2 + \eps\|\vx_{k+1}-\vx^*\|^2 +  \frac{E^2 + \zeta^2}{2\eps} \\
        & + D_{h}(\vx_{k+1};\vx_k ) - \frac{1}{2\eta}\|\vx_{k} - \vx_{k+1}\|^2 . 
    \end{align*}
    Which, using the definition of $\Gamma_k$ (\cref{app:eq:def_Lyapunov}), yields
    \begin{align}
         \left(1 - \frac{\alpha \mu}{2 \eps}\right)\left(f(\vx_{k+1}) - f(\vx*)\right) &\le \Gamma_k - \Gamma_{k+1} - \left(1 - \frac{\alpha \mu}{2 \eps}\right) D_f(\vx^*;\vx_{k+1}) \notag\\
    & + \frac{c}{2\eps}\|\vx_{k+1} - \vx_k\|^2 + \eps\|\vx_{k+1}-\vx^*\|^2 +  \frac{E^2 + \zeta^2}{2\eps} \notag\\
    & + D_{h}(\vx_{k+1};\vx_k ) - \frac{1}{2\eta}\|\vx_{k} - \vx_{k+1}\|^2. \label{app:eq:decomposition_withlyap_1}
    \end{align}
    
    
    We now use the regularity property of $f$ and $h$ to derive the following inequalities: 
    \begin{align}
        D_{f}(\vx^*;\vx_{k+1}) &\ge \frac{\mu}{2}\|\vx_{k+1}-\vx^*\|^2, &&\text{since $f$ is $\mu$-strongly convex}\label{app:eq:usage_strong_cvx_f}\\
        |D_{h}(\vx^*;\vx_{k+1})| &\le \frac{L_h}{2}\|\vx_{k+1}-\vx^*\|^2,&&\text{since $h$ is $L_h$-smooth. }\label{app:eq:usage_smooth_h}
    \end{align}
    And, using \cref{app:eq:usage_smooth_h}, we have, for $\eta^{-1} \ge L_h$,
    \[
    0 \le \frac{1}{2\eta}\|\vx_{k+1} - \vx^*\|^2 - D_h(\vx^*,\vx_{k+1}) \le \frac{\eta^{-1} + L_h}{2}\|\vx_{k+1} - \vx^*\|^2.
    \]
    Which yields, considering that $D_f(\vx^*;\vx_{k+1}) + D_f(\vx_{k+1},\vx^*)=\langle \nabla f(\vx_{k+1}), \vx_{k+1}-\vx^*\rangle$, and using \cref{app:eq:usage_strong_cvx_f},
    \begin{align}
        0 \le \Gamma_{k+1} &\le \frac{\eta^{-1} + L_h}{2}\|\vx_{k+1} - \vx^*\|^2 + \frac{\alpha \mu}{2 \eps}[\langle \nabla f(\vx_{k+1}), \vx_{k+1}-\vx*\rangle]\notag\\
        0 \le \Gamma_{k+1}&\le \left(\frac{\eta^{-1} + L_h}{\mu} + \frac{\alpha \mu}{2 \eps}\right) D_f(\vx^*;\vx_{k+1}) + \frac{\alpha \mu}{2 \eps}(f(\vx_{k+1}) - f^*)\label{app:eq:usage_regularity_on_Gamma} \ .
    \end{align}
    
Combining \cref{app:eq:usage_strong_cvx_f} and \cref{app:eq:usage_regularity_on_Gamma} yields
\begin{equation}
\label{app:eq:lowerbound_Df}
D_f(\vx^*;\vx_{k+1}) \ge  \frac{\mu}{4}\|\vx_{k+1}-\vx^*\|^2 +\frac{1}{2}\left(\frac{\eta^{-1} + L_h}{\mu} + \frac{\alpha \mu}{2 \eps}\right)^{-1}  \left(\Gamma_{k+1} - \frac{\alpha \mu}{2 \eps} (f(\vx_{k+1})-f(\vx^*)\right).
\end{equation}

We can now combine \cref{app:eq:lowerbound_Df} and \cref{app:eq:decomposition_withlyap_1}.
\begin{align*}
    \left(1 - \frac{\alpha \mu}{2 \eps}\right)\left(f(\vx_{k+1}) - f(\vx*)\right) &\le \Gamma_k - \Gamma_{k+1} \\
    &- \left(1 - \frac{\alpha \mu}{2 \eps}\right) \left[\frac{\mu}{4}\|\vx_{k+1}-\vx^*\|^2 +\frac{1}{2}\left(\frac{\eta^{-1} + L_h}{\mu} + \frac{\alpha \mu}{2 \eps}\right)^{-1}  \left(\Gamma_{k+1} - \frac{\alpha \mu}{2 \eps} (f(\vx_{k+1})-f(\vx^*)\right)\right] \notag\\
    & + \frac{c}{2\eps}\|\vx_{k+1} - \vx_k\|^2 + \eps\|\vx_{k+1}-\vx^*\|^2 +  \frac{E^2 + \zeta^2}{2\eps} \notag\\
    & + D_{h}(\vx_{k+1};\vx_k ) - \frac{1}{2\eta}\|\vx_{k} - \vx_{k+1}\|^2. 
\end{align*}
For simplicity, we denote $\rho = \frac{\mu}{\eta^{-1}  + L_h}$ and $\tau = \frac{\alpha \mu}{2 \eps}$, which yields 
\begin{align*}
    \left(1 - \tau\right)\left(f(\vx_{k+1}) - f(\vx*)\right) &\le \Gamma_k - \Gamma_{k+1} \\
    &- \left(1 - \tau\right) \left[\frac{\mu}{4}\|\vx_{k+1}-\vx^*\|^2 +\frac{1}{2}\left(\rho + \tau\right)^{-1}  \left(\Gamma_{k+1} - \tau (f(\vx_{k+1})-f(\vx^*)\right)\right] \notag\\
    & + \frac{c}{2\eps}\|\vx_{k+1} - \vx_k\|^2 + \eps\|\vx_{k+1}-\vx^*\|^2 +  \frac{E^2 + \zeta^2}{2\eps} \notag\\
    & + D_{h}(\vx_{k+1};\vx_k ) - \frac{1}{2\eta}\|\vx_{k} - \vx_{k+1}\|^2. 
\end{align*}
Re-arranging it yields
\begin{align*}
        \left(1 - \tau\right)\left(1 - \frac{\tau}{2(\rho^{-1} + \tau)}\right)&(f(\vx_{k+1}) - f(\vx*)) \\
        &\le \Gamma_k - \left(1 +  \frac{1-\tau}{2(\rho^{-1} + \tau)}\right)\Gamma_{k+1} + \frac{\zeta^2}{2 \eps} + \frac{E^2}{2\eps}\\
        &  \left( \frac{c}{2\eps} + \frac{L_h}{2} - \frac{1}{2\eta} \right)\|\vx_{k} - \vx_{k+1}\|^2 + \left( \eps - \left(1 - \tau\right) \frac{\mu}{4} \right)\|\vx_{k+1}-\vx^*\|^2.
\end{align*}
We take the largest $\eps$ such that  $(\eps - (1-\frac{\alpha \mu}{2\eps})\frac{\mu}{4}) \le 0$:
\[
\eps^2 - \frac{\mu}{4} \eps +\frac{\alpha \mu^2}{8} = 0 \iff \eps = \frac{\mu \pm \sqrt{\mu^2(1-8\alpha^2)}}{8} \qquad \text{when} \quad \alpha < 1/8.
\]
We thus take $\eps = \frac{ \mu}{8}(1 + \sqrt{1 - 8 \alpha}) \ge \frac{\mu}{8}$.

Moreover, the coefficient in front of $\|\vx_k - \vx_{k+1}\|^2$ is non-positive if $\eta^{-1} \ge L_h + c/\eps$, thus for $\eta\le  1/(L_h + \frac{8c}{\mu})$.

This yields 
\begin{align*}
        \left(1 - \tau\right)\left(1 - \frac{\tau}{2(\rho^{-1} + \tau)}\right)&(f(\vx_{k+1}) - f(\vx^*)) \\
        &\le \Gamma_k - \left(1 +  \frac{1-\tau}{2(\rho^{-1} + \tau)}\right)\Gamma_{k+1} +  \frac{4(\zeta^2 +E^2)}{\mu}.
\end{align*}

Which can be additionally simplified, by noting that $\tau = \frac{\alpha \mu}{2\eps} \le 4 \alpha \le \frac{1}{2}$, and \(\frac{\tau}{2(\rho^{-1} + \tau)} \le \frac{1}{2}\) and $\rho = \frac{\eta \mu}{1+\eta L_h} \ge \frac{\eta\mu}{2}$, into

\begin{align*}
        \frac{1}{4}(f(\vx_{k+1}) - f(\vx^*)) 
        \le \Gamma_k - \left(1 +  \frac{\eta \mu}{8}\right)\Gamma_{k+1} + \frac{4\zeta^2}{\mu} + \frac{4E^2}{\mu}.
\end{align*}
Which is the main result.

\end{proof}


We can now prove \cref{app:thm:proxy_prox_inexact}.

\begin{proof}

Denoting $\beta_k := \left(1 + \frac{\eta \mu}{8}\right)^k$ and $B_k := \sum_{k=0}^K \beta_k$ and considering
\[
\hat{\vx}_{K} := \frac{1}{B_K}\sum_{k=1}^K \beta_k \vx_k,
\]
allows to write, by convexity of $f$,
\begin{align*}
    f(\hat{\vx}_K) - f^* &\le \frac{1}{B_k}\sum_{k=1}^K \beta_k (f(\vx_k) - f^*)\\
    &\le 4\left(\frac{\Gamma_{0}}{B_K} + \frac{4\zeta^2}{\mu} + \frac{4E^2}{\mu}\right). &&\text{using \cref{app:lemma:proof_proxyprox_3}}
\end{align*}

Now, using $ \tau \le \frac{1}{2}$ and the $L_f$-smoothness of $f$, we have $\Gamma_{0} \le \frac{\eta^{-1} + L_f}{2}\|\vx_0 - \vx^*\|^2 \le \frac{1}{\eta}\|\vx_0 - \vx^*\|^2$, and
\begin{align*}
    B_K &= \frac{(1+\frac{\eta \mu}{8})^K - 1}{\eta \mu /8} = \frac{(1+\frac{\eta \mu}{8})^{K-1}\frac{\eta \mu}{8} + (1+\frac{\eta \mu}{8})^{K-1} - 1}{\eta \mu /8}\\
    & \le \left(1 + \frac{\eta \mu }{8}\right)^{K-1} \ .
\end{align*}
We have 
\begin{align*}
f(\hat\vx_{K}) - f(\vx^*) &\le 4\left(1 + \frac{\eta \mu}{8}\right)^{1-K}\Gamma_0 + 16\frac{\zeta^2 + E^2}{\mu} \\
f(\hat\vx_{K}) - f(\vx^*) &\le \left(1 + \frac{\eta \mu}{8}\right)^{1-K}2(\eta^{-1} + L_f)\|\vx_0 - \vx^*\|^2 + 16\frac{\zeta^2 + E^2}{\mu}.
\end{align*}
Which concludes the proof.
\end{proof}



\subsubsection{Proof of \cref{prop:similarity_to_GB_heter}}


\label{app:proof:similarity_to_GB_heter}
Let's write the differences of gradients as an integral.
\begin{align*}
    \nabla \cL_i(\vx) - \nabla \cL_{\cH}(\vx) &= \nabla \cL_i(\vx^*) + \int_{0}^1 \nabla^2\left(\cL_i - \cL_{\cH}\right)(\vx^* + t(\vx-\vx^*) ) \cdot (\vx-\vx^*)dt.
\end{align*}
 We now use Young's inequality, then apply Jensen's inequality. This allows us to leverage the similarity \cref{asmpt:Pariwise_Hessian_Similarity}, which, along with the convexity of both $\cL_i$ and $\cL_{\cH}$, concludes the proof.
\begin{align*}
    \left\|\nabla \cL_i(\vx) - \nabla \cL_{\cH}(\vx)\right\|^2 &\le 2\left\|\nabla \cL_i(\vx^*)\right\|^2 + 2\left\|\int_{0}^1 \nabla^2\left(\cL_i - \cL_{\cH}\right)(\vx^* + t(\vx-\vx^*) ) \cdot (\vx-\vx^*)dt\right\|^2\\
    &\le 2\left\|\nabla \cL_i(\vx^*)\right\|^2 + 2\int_{0}^1 (\vx-\vx^*)^T (\nabla^2\left(\cL_i - \cL_{\cH}\right)(\vx^* + t(\vx-\vx^*) ))^2 \cdot (\vx-\vx^*)dt\\
    &\le 2\left\|\nabla \cL_i(\vx^*)\right\|^2 + 2\Delta \int_{0}^1 |(\vx-\vx^*)^T \nabla^2\left(\cL_i - \cL_{\cH}\right)(\vx^* + t(\vx-\vx^*) ) \cdot (\vx-\vx^*)|dt\\
    &\le 2\left\|\nabla \cL_i(\vx^*)\right\|^2 + 2\Delta \int_{0}^1 (\vx-\vx^*)^T \nabla^2\left(\cL_i + \cL_{\cH}\right)(\vx^* + t(\vx-\vx^*) ) \cdot (\vx-\vx^*)dt\\
    &= 2\left\|\nabla \cL_i(\vx^*)\right\|^2 + 2\Delta \langle \nabla \cL_i(\vx) + \nabla \cL_{\cH}(\vx) - \nabla \cL_i(\vx^*), \vx-\vx^*\rangle.
\end{align*}
Averaging it on all $i \in \cH$, and using that $\nabla \cL_{\cH}(\vx^*) \triangleq 0$, yields
\[
\frac{1}{|\cH|}\sum_{i \in \cH}\left\|\nabla \cL_i(\vx) - \nabla \cL_{\cH}(\vx)\right\|^2 \le \frac{2}{|\cH|}\sum_{i \in \cH}\|\nabla \cL_i(\vx^*)\|^2 + 4\Delta \langle \nabla \cL_{\cH}(\vx),\vx-\vx^*\rangle. 
\]
Using \cref{app:eq:cocoaxivity_strong_convexity}, written here as $\langle \nabla \cL_{\cH}(\vx), \vx-\vx^*\rangle \le  \frac{1}{\mu}\|\nabla \cL(\vx)\|^2$, gives the result. 

\subsection{Simplification of \cref{alg:fast_gradient_method}}
\label{app:sec:algorithm_simplification}
In \citet{devolder2014first}, \cref{alg:fast_gradient_method} writes as follows:
Given $(\gamma_k)_k$, $(\Gamma_k)_k$, $(\tau_k)_k$ satisfying:
\begin{align*}
    \Gamma_k := \sum_{i=0}^k\gamma_i\\
    L + \mu \Gamma_k = \frac{L \gamma_{k+1}^2}{\Gamma_{k+1}}\\
    \tau_k := \frac{\gamma_{k+1}}{\Gamma_{k+1}} = \frac{L + \mu\Gamma_k}{\Gamma_{k+1}}, 
\end{align*}

\cref{alg:fast_gradient_method} writes, given a $(\tL,\tmu,\delta)$ inexact gradient oracle:

\begin{algorithm}[h]
  \caption{Inexact Oracle Fast Gradient Method}
  \label{app:alg:fast_gradient_method_devolder}
  \begin{algorithmic}
    \STATE {\bfseries Input:} $(\tilde\nabla \cL)$, $\eta$, $\beta$, $\gamma_0 = 1$, $(\tau_k)_k$, $(\gamma_k)$, $\vx_0$
     \FOR{$k=1$ {\bfseries to} $K$}
     \STATE Sample $\tilde\nabla \cL_{\cH}(\vx_k)$
    \STATE Compute $\vy_{k} = \argmin_{\vy}\left\{\langle \vx -\vx_k, \tilde{\nabla}\cL_{\cH}(\vx_k)\rangle + \frac{\tL}{2}\|\vx - \vx_k\|^2\right\}$ 
    \STATE Set $\phi_k(\vx) = \frac{\tL}{2}\|\vx-\vx_0\|^2 + \sum_{i=0}^k \gamma_i[\langle\tilde{\nabla}\cL_{\cH}(\vx_i), \vx - \vx_i\rangle + \frac{\tmu}{2}\|\vx - \vx_i\|^2$
    \STATE Compute $\vz_k = \argmin_{\R^d}\phi_k(\vx)$
    \STATE Compute $\vx_{k+1} = (1-\tau_k)\vy_k + \tau_k \vz_k$
    \ENDFOR
  \end{algorithmic}
\end{algorithm}
We now show write \cref{app:alg:fast_gradient_method_devolder} without $\argmin$ operations.

First, we have that 
\[
\tilde{\nabla}\cL_{\cH}(\vx_k) + \tL(\vy_k - \vx_k) = 0 \iff \vy_k = \vx_k - \frac{1}{\tL}\tilde{\nabla}\cL_{\cH}(\vx_k).
\]
Then the definition of $\vz_k$ is
\begin{align*}
    \tL(\vz_k - \vx_0) + \sum_{i=0}^k \gamma_i[\tilde{\nabla}\cL_{\cH}(\vx_i) + \tmu (\vz_k - \vx_i)] = 0 \iff (\tL +\mu \Gamma_i) \vz_k = \tL\vx_0 + \sum_{i=0}^k \gamma_i(\tmu \vx_i - \tilde{\nabla}\cL_{\cH}(\vx_k)).
\end{align*}
This yields 
\begin{align*}
    (\tL +\mu \Gamma_k) \vz_k &= (\tL +\mu \Gamma_{k-1}) \vz_{k-1} + \gamma_k(\tmu \vx_k - \tilde{\nabla}\cL_{\cH}(\vx_k))\\
    &= (\tL +\mu \Gamma_{k-1}) \vz_{k-1} + \gamma_k( \tL \vy_k - (\tL-\tmu) \vx_k).
\end{align*}
Plugging-in $\vz_{k-1} = \frac{\vx_k - (1-\tau_{k-1})\vy_{k-1}}{\tau_{k-1}}$ yields
\begin{equation}
\label{app:eq:identity1}
\tau_k\vz_k = \frac{\tau_k}{\tau_{k-1}}\frac{\tL +\mu \Gamma_{k-1}}{\tL +\mu \Gamma_{k}} (\vx_k - (1-\tau_{k-1})\vy_{k-1}) + \frac{\tL\gamma_k \tau_k}{\tL +\mu \Gamma_{k}}(  \vy_k - (1-\frac{\tmu}{\tL}) \vx_k).
\end{equation}

Grouping the terms in front of $\vx_k$ yields:
\begin{align*}
    \frac{1}{\tL + \tmu \Gamma_k}\left( \frac{\tau_k}{\tau_{k-1}} (L\gamma_{k}\tau_{k-1}) - \tL \gamma_k \tau_k(1-\frac{\tmu}{\tL})\right)= -\frac{\tmu}{\tL}\frac{\tL\gamma_k \tau_k}{\tL + \tmu\Gamma_k}.
\end{align*}

And, since, 
\begin{equation}
\label{app:eq:identity3}
    \frac{\tL\gamma_k \tau_k}{\tL +\mu \Gamma_{k}} = \frac{\tL\gamma_k \tau_k}{\tL\tau_{k}\gamma_{k+1}} = \frac{\gamma_k}{\gamma_{k+1}}.
\end{equation}

Thus \cref{app:eq:identity1} can be simplified as 
\begin{equation}
\label{app:eq:identity4}
    \tau_k\vz_k =  \frac{\gamma_k}{\gamma_{k+1}}\frac{\tmu}{\tL}\vx_k - \frac{\tau_k}{\tau_{k-1}}\frac{\tL +\mu \Gamma_{k-1}}{\tL +\mu \Gamma_{k}}  (1-\tau_{k-1})\vy_{k-1} + \frac{\gamma_k}{\gamma_{k+1}}\vy_k .
\end{equation}

Eventually, considering that
\begin{align*}
    \frac{\tau_k}{\tau_{k-1}}\frac{\tL +\mu \Gamma_{k-1}}{\tL +\mu \Gamma_{k}}  (1-\tau_{k-1}) 
    &=\frac{\tau_k}{\tau_{k-1}}\frac{L\tau_{k-1}\gamma_k}{L\tau_{k}\gamma_{k+1}}  \frac{\Gamma_k - \gamma_k}{\Gamma_k}\\
    &= \frac{\gamma_k \Gamma_{k-1}}{\gamma_{k+1}\Gamma_k},
\end{align*}
we have
\[
\tau_k \vz_k = \frac{\gamma_k}{\gamma_{k+1}}(\frac{\tmu}{\tL}\vx_k + \vy_k - \frac{\Gamma_{k-1}}{\Gamma_k}\vy_{k-1}).
\]
Plugging this in $\vx_{k+1} = (1-\tau_k)\vy_k + \tau_k \vz_k$ yields

\begin{align*}
    \vx_{k+1} &= \frac{\gamma_k}{\gamma_{k+1}}( \frac{\tmu}{\tL} \vx_k + \vy_k) + \frac{\Gamma_{k-1}}{\Gamma_k}\left( \vy_k - \frac{\gamma_k}{\gamma_{k+1}} \vy_{k-1}\right)\\
    \vx_{k+1}&=\vy_k + \frac{\Gamma_{k-1}}{\Gamma_k}\frac{\gamma_k}{\gamma_{k+1}}(\vy_k - \vy_{k-1}) - \frac{\Gamma_{k-1}}{\Gamma_k}\vy_k + \frac{\gamma_k}{\gamma_{k+1}}\frac{\tmu}{\tL}\vx_k .
\end{align*}
 Which, using $\tL = 2L$ and $\tmu = \mu/2$, gives the following algorithm:
\begin{algorithm}[h]
  \caption{Byzantine-Resilient Fast Gradient Method}
  \begin{algorithmic}
    \STATE {\bfseries Input:} $F$, $\eta$, $\beta$, $\gamma_0 = 1$, $(\tau_k)_k$, $(\gamma_k)$, $\vx_0$
     \FOR{$k=1$ {\bfseries to} $K$}
     \STATE Sample $\tilde\nabla \cL_{\cH}(\vx_k)$ using \cref{alg:reduction}
    \STATE Compute $\vy_{k} = \vx_k - \frac{1}{2L}\tilde{\nabla}\cL_{\cH}(\vx_k)$ 
    \STATE Compute $\vx_{k+1}=\vy_k + \frac{\Gamma_{k-1}}{\Gamma_k}\frac{\gamma_k}{\gamma_{k+1}}(\vy_k - \vy_{k-1}) - \frac{\Gamma_{k-1}}{\Gamma_k}\vy_k + \frac{\gamma_k}{\gamma_{k+1}}\frac{\mu}{4L}\vx_k$
    \ENDFOR
  \end{algorithmic}
\end{algorithm}
\section{Experiments}
\label{app:experiments}


\begin{figure}[h]
\centering
\includegraphics[width=0.8\textwidth]{mnist_logreg_mnist_n_21_f_1_dirichlet_niid_5.0_wd_0.001_TrMean_NNM_comparison_test_accuracy.pdf}
\caption{Distributed Logistic Regression, with $20$ honest and $1$ Byzantine client in the mildly heterogeneous setting ($\beta=5$), with a $l2$-regularization $\lambda = 10^{-3}$. Test Accuracy under A Little Is Enough (ALIE) attack.}
\label{plot:logistic_heterogeneous_accuracy}
\end{figure}


\begin{figure*}[h]
\centering
\includegraphics[width=0.9\textwidth]{mnist_logreg_mnist_n_21_f_1_dirichlet_niid_1.0_wd_0.001_TrMean_NNM_comparison_with_ref_test_accuracy.pdf}
\caption{Impact of the learning rate $\eta$ on PIGS, in a highly heterogeneous setting $(\beta=1)$ with $20$ honest clients and either with $1$ Byzantine performing IPM or ALIE attack, or without Byzantine client. Penalization $\mu = 10^{-3}$. Test Accuracy.}
\end{figure*}


\begin{figure*}[h]
\centering
\includegraphics[width=0.9\textwidth]{mnist_logreg_mnist_n_25_f_5_iid_None_wd_0.001_TrMean_NNM_comparison_with_ref_test_accuracy.pdf}
\caption{Impact of the learning rate $\eta$ on PIGS,i.i.d. setting with $20$ honest clients and either with $5$ Byzantine performing IPM or ALIE attack, or without Byzantine client. Penalization $\mu = 10^{-3}$. Test Accuracy.}
\end{figure*}





\end{document}
